% environment.tex — ループ計算環境の構築記録（2026年10月9日）
% コンパイル: lualatex environment.tex（参照解決のため2回）
% 付録Cは files/ 以下のファイルを \VerbatimInput で取り込む。
\documentclass[a4paper,10pt]{ltjsarticle}
\usepackage[margin=20mm]{geometry}
\usepackage[haranoaji,deluxe]{luatexja-preset}
\ltjsetparameter{jacharrange={-3}}% 一般句読点（– — ‘ ’ など）は欧文フォントで組む
\usepackage{amsmath,amssymb}
\usepackage{booktabs,tabularx,array}
\usepackage[table]{xcolor}
\usepackage{tikz}
\usetikzlibrary{positioning,arrows.meta,calc,fit,backgrounds}
\usepackage[font=small,labelfont=bf]{caption}
\usepackage{fancyvrb,fvextra}
\usepackage{xurl}
\usepackage{hyperref}
\hypersetup{colorlinks=true,linkcolor=blue!50!black,citecolor=green!40!black,urlcolor=blue!60!black,
  pdftitle={ループ計算環境の構築記録},pdfauthor={}}

% 色（Mathematica 依存）
\definecolor{cMf}{RGB}{253,229,204}\definecolor{cMd}{RGB}{204,112,26}  % Mathematica 必須
\definecolor{cOf}{RGB}{214,232,247}\definecolor{cOd}{RGB}{46,99,158}   % Mathematica 不要
\definecolor{cPf}{RGB}{228,228,228}\definecolor{cPd}{RGB}{110,110,110} % 補助
\definecolor{cGf}{RGB}{226,243,229}\definecolor{cGd}{RGB}{52,128,72}   % 成功・一致
\definecolor{cRf}{RGB}{252,226,226}\definecolor{cRd}{RGB}{190,30,30}   % 失敗
\definecolor{cQf}{RGB}{255,246,214}\definecolor{cQd}{RGB}{176,140,20}  % 手順

\tikzset{
  tool/.style={draw,rounded corners=1.5pt,minimum height=0.48cm,text width=4.6cm,align=flush center,
    inner sep=1.6pt,line width=0.5pt,font=\footnotesize},
  tM/.style={tool,fill=cMf,draw=cMd},
  tO/.style={tool,fill=cOf,draw=cOd},
  tN/.style={tool,fill=white,draw=cPd},
  flow/.style={-{Stealth[length=2mm]},thick,black!70},
  panel/.style={fill=black!3,draw=black!22,rounded corners=3pt},
  ptitle/.style={font=\footnotesize\bfseries,anchor=north,align=flush center,inner sep=2pt},
}

\newcommand{\MGaMC}{MG5\_aMC@NLO}
\newcommand{\bq}{\textasciigrave}
\newcommand{\pth}[1]{\texttt{\textasciitilde/}\nolinkurl{#1}}
\renewcommand{\jsParagraphMark}{\textcolor{black!65}{\rule[-0.02\zw]{0.62\zw}{0.62\zw}}\hspace{0.35\zw}}
\fvset{fontsize=\footnotesize,frame=lines,framesep=2mm,breaklines=true,breakanywhere=true}

\title{ループ計算環境の構築記録\\[2pt]{\large \texttt{\textasciitilde/hep} に組んだ1ループ計算スタック（Apple M3・Mathematica 14.2）}}
\author{}
\date{2026年10月9日（熱的な道具の追加は10月11日）}

\begin{document}
\maketitle

\begin{abstract}
本稿は、Feynman 図の1ループ計算に使うソフトウェア環境を、この Mac（Apple M3、macOS 26.2、Mathematica 14.2.1）に組み直した記録である。旧環境には FeynCalc 10.1、FeynArts、Package-X、FeynRules、LieART、FormCalc、\MGaMC{} 3.6.2 があり、そのインストールとアーカイブは計22項目に散らばっていた。自作の模型ファイル570個を退避して元と照合したのち、これらをすべてゴミ箱へ移し、同じリポジトリのレビューが推奨した1ループの構成を導入した。すなわち FeynCalc 10.2.1 + FeynHelpers 2、Package-X 2.1.1、LoopTools 2.16、COLLIER 1.2.10、FORM 5.0.2、FormCalc 9.10、FeynRules 2.3.49、\MGaMC{} 3.8.1 である。導入の途中で、この Mac に固有の障害が四つ見つかった。gfortran が macOS の SDK を見つけられないこと、Mathematica 付属の MathLink コンパイラ \texttt{mcc} が Apple Silicon に対応していないこと、Mathematica 14.1 で入った組み込み記号 \texttt{MatrixSymbol} が FeynRules の同名記号を隠すこと、ブラウザで取得した実行ファイルが macOS に削除されることである。それぞれ、SDKROOT の設定、代替の \texttt{mcc}、FeynRules.m への局所パッチ、curl による取得で回避した。検証では三つの点を確かめた。第一に、$B_0(2;1,1)=2-\pi/2$ を LoopTools・COLLIER・Package-X・FeynHelpers の各経路で再現した（倍精度の全桁、COLLIER は出力した12桁）。第二に、MadGraph の $e^+e^-\to\mu^+\mu^-$ の断面積 $0.10437\pm0.00018$~pb が、$\gamma$ と $Z$ の交換を含む解析式の値 $0.10426$~pb と $0.6\sigma$ 以内で一致した。第三に、FeynRules が出力した標準模型の UFO が、MadGraph 同梱の標準模型の153頂点をすべて含むことを確かめた。その後、$\phi^3$ 模型での検算で、この環境の LoopTools が閾値より下で外運動量がゼロでない $C_0$ を誤って返し、その誤りがテンソル係数 $C_{00}$ にも及ぶことが分かった。$C_0$ と $C_{00}$ には Package-X と COLLIER を使う。4点関数 $D_0$ は、外線と内線の質量がすべて等しいとき、ライブラリごとに異なる点で値を返さない。さらに MadLoop との比較で、NLOCT の不具合（FeynRules の \texttt{Core/Loop.m} に2つ目の局所パッチを当てた）と、MadGraph の1ループ計算が QCD と QED の摂動しか受け付けないことが分かった。全断面積の計算では、MadGraph の木の $2\to3$ の断面積が、高いエネルギーで誤差を小さく見積もり、自前の積分より $0.8\%$ 低く出た。$\phi\phi\to\phi\phi$ を高エネルギーまで計算すると、COLLIER の箱図が $\sqrt s\gg m$ で精度を失った。既定の方式 COLI では虚部が、もう一方の DD では $\sqrt s\gtrsim10^3m$ で値全体が誤る。LoopTools と、Package-X の解析式を30桁で評価した値は正しかった（第~\ref{sec:issues} 節）。計算のスクリプトは、まとめて実行できるうえ、Mathematica のノートブックや VS Code でセルごとにも実行できる形にした。これで途中の式や図を確かめられる。$\phi^3$ 模型の25本すべてをこの形にし、結果が変わらないことを確かめた。評価する Feynman 図は、評価の前のセルで表示する。そのための道具は \nolinkurl{environment/interactive/} にある（第~\ref{sec:usage} 節）。全検証は \nolinkurl{environment/tests/run_all.sh} で再実行できる。2026年10月11日には、有限温度の相転移の検算のために、熱的な道具を加えた。DRalgo 1.5.0（GroupMath 1.1.3 とともに）、Matchete 0.6.0 の Matchotter、PT2GWFinder 1.1.0 と FindBounce 1.1.0、CosmoTransitions 2.0.7 である。DRalgo と Matchotter は、$\phi^3+\phi^4$ 理論の有限温度のノートの温度ポテンシャルの係数と記号で一致した。CosmoTransitions は bounce の作用と転移の温度を、ノートの自前の計算と相対 $10^{-7}$ 以内で再現した。PT2GWFinder は、臨界温度と相転移の強さを相対 $10^{-6}$ 以内、核生成温度を $2.4\times10^{-4}$ 以内で再現した。さらに PhaseTracer2 2.2.2、BSMPT 3.2.0、TransitionListener 2.1.0 も加え、臨界温度、作用、強さ、$\beta/H$ が、核生成の条件や $g_*$ などの取り決めの違いを除いて再現されることを確かめた。FindBounce は、経路の端を深い真の極小に置くと、点を増やすほど正しい値から離れた。Matchotter の対数の尺度は通常の $\overline{\mathrm{MS}}$ の尺度と $e^{\gamma_E}$ 倍違い、これは論文の付録 B の規格化と一致した。最後に WallGo 1.1.2 を加え、局所熱平衡の壁の速さを求めた。TransitionListener の値との差は、TransitionListener が音速を場を固めた温度微分で求めることから生じると分かった（第~\ref{sec:thermal} 節）。10月12日には、2ループの道具として Kira 3.1（Fermat 7.9b とともに）、FIRE 7.1、AMFlow 2.0、pySecDec 1.6.6、feyntrop 1.1.0、TSIL 1.46 を加えた。$\phi^3$ 理論の $d=4$ の2ループの自己エネルギーで、Kira と FIRE の IBP 簡約は記号として一致し、マスター積分は四つの数値法で一致し、有限部分は文献の式と相対 $8\times10^{-14}$ 以内で一致した。Apple clang 16 で \texttt{-O3} で構築した Kira は、すべての積分を0に簡約した（第~\ref{sec:twoloop} 節）。
\end{abstract}

\section{はじめに}\label{sec:intro}

本稿の問いは次の一つである。\textbf{この Mac で1ループ計算を正しく、かつ再現可能に行うために、どのような環境を組み、それが正しく動くことをどう確かめたか。}環境を組み直した理由は、旧環境に五つの問題があったことである。FeynCalc は 10.1.0 で、外部ツールとの連携を担う FeynHelpers が入っていなかった。LoopTools はビルドされておらず、Package-X は難読化された旧配布版だった。FormCalc に付属する FORM の実行ファイルは、検証のために起動した直後に macOS によって削除された。FeynRules は Mathematica 14 で読み込むたびにエラーを出した。さらに、同じツールのコピーが Library・Documents・Downloads・Dropbox に散らばり、どれが使われているのか分からなかった。そこで同じリポジトリのレビュー（\nolinkurl{loop_tools_review/loop_tools_review.pdf}）が推奨した構成のうち、1ループの中核部分を導入することにした。2ループ用のツールと、EFT・暗黒物質用のツールは、本稿の範囲外である。熱的ループと相転移の道具は、のちに有限温度の検算のために加えた（第~\ref{sec:thermal} 節）。

\paragraph{設定。}対象とする計算機と、ビルドに使った共通の基盤を表~\ref{tab:machine} に示す。Mathematica は 14.2.1（\texttt{/Applications/Wolfram 2.app}）と 14.1.0（\texttt{/Applications/Wolfram.app}）の二つが入っている。両者は同じ \texttt{\$UserBaseDirectory}（\pth{Library/Wolfram}）を参照するので、パッケージはどちらからも見える。ビルドと検証には 14.2.1 だけを使った。

\begin{table}[tbp]
\centering\small
\caption{計算機と共通の基盤（2026年10月9日時点）。}
\label{tab:machine}
\begin{tabularx}{\linewidth}{>{\raggedright\arraybackslash}p{4.0cm}X}
\toprule
\textbf{項目} & \textbf{内容} \\
\midrule
計算機 & Apple M3（8コア）、メモリ 16 GB \\
OS & macOS 26.2。Command Line Tools の SDK は 15.2 \\
Mathematica & Wolfram 14.2.1（ビルドと検証に使用）、14.1.0 \\
コンパイラ & Homebrew gcc 15.2.0（gfortran）、Apple clang 16.0.0、cmake 3.31.1 \\
数値ライブラリ & GMP 6.3.0、MPFR 4.2.2、FLINT 3.6.0、zstd 1.5.7（いずれも Homebrew） \\
Python & miniforge（conda 24.11.3）の環境 \texttt{qft}：Python 3.12.15、numpy 2.5.3、scipy 1.18.1、six 1.17.0。図と高精度の検算のため、のちに matplotlib 3.11.2 と mpmath 1.4.1 を加えた \\
\bottomrule
\end{tabularx}
\end{table}

\paragraph{前提。}本稿は次の前提に立つ。
\begin{enumerate}
\item 版と日付は2026年10月9日時点のものである。熱的な道具（第~\ref{sec:thermal} 節）は10月11日、2ループの道具（第~\ref{sec:twoloop} 節）は10月12日時点である。
\item ソフトウェアは公式の配布元から \texttt{curl} で取得した。ソースから構築できるものはこの Mac で構築した。
\item 「正しく動く」かどうかは、既知の厳密値、独立な別実装、または独立な解析計算との数値の一致で判定する。
\item パスの \texttt{\textasciitilde} はホームディレクトリ \nolinkurl{/Users/koichiro} を表す。
\end{enumerate}
本文で使う用語と記号を表~\ref{tab:terms} にまとめる。

\begin{table}[tbp]
\centering\small
\caption{本稿で使う用語と記号。}
\label{tab:terms}
\begin{tabularx}{\linewidth}{>{\raggedright\arraybackslash}p{3.6cm}X}
\toprule
\textbf{用語・記号} & \textbf{意味} \\
\midrule
\texttt{\textasciitilde/hep} & 今回の構成で、Mathematica パッケージ以外を置いた場所（ソース、ビルド結果、MadGraph） \\
Applications & \pth{Library/Wolfram/Applications}。Mathematica がパッケージを探す場所 \\
MathLink, \texttt{mcc} & Mathematica と外部プログラムをつなぐ通信規格と、そのためのコンパイラ・スクリプト \\
UFO & 模型を MadGraph などへ渡す Python 形式（Universal FeynRules Output） \\
HEPTools & MadGraph が1ループ計算に使う外部ライブラリ群（Ninja、OneLOop、COLLIER） \\
PaVe, $A_0, B_0, C_0$ & Passarino--Veltman のスカラー関数。規格化は LoopTools に合わせ、$\mu^2=1$、$\Delta=0$ とする（\ref{app:derivations}） \\
$B_0(s;m_0,m_1)$ & 外運動量の2乗 $s$、内線質量 $m_0,m_1$ の2点関数。本稿では $B_0(2;1,1)$ を検算に使う \\
SDKROOT & gfortran や clang が使う macOS SDK の場所を示す環境変数 \\
quarantine & macOS がダウンロードしたファイルに付ける拡張属性 \texttt{com.apple.quarantine} \\
\bottomrule
\end{tabularx}
\end{table}

本稿の構成は次の通りである。第~\ref{sec:migration} 節で旧環境の撤去と自作模型の退避を、第~\ref{sec:layout} 節で導入した構成を述べる。第~\ref{sec:issues} 節では、この Mac に固有の障害とその対処を、第~\ref{sec:tests} 節では検証の方法と結果を示す。第~\ref{sec:usage} 節に使い方を、第~\ref{sec:thermal} 節と第~\ref{sec:twoloop} 節に後から加えた熱的な道具と2ループの道具を、第~\ref{sec:discussion} 節に検証していない点と今後の課題をまとめる。付録には、ファイル配置、検算に使った解析式の導出、変更・作成したファイルの全文を載せる。

\section{旧環境の撤去と自作模型の退避}\label{sec:migration}

撤去の手順を図~\ref{fig:migration} に示す。最も注意したのは、ツールのフォルダの中に自作のファイルが紛れていないかである。FeynRules と FeynArts は、利用者の模型を自分の \texttt{Models} フォルダに置く使い方を前提にしており、MadGraph も計算結果をインストール先に書き出すことがある。そこで、展開済みのフォルダを元の tarball のファイル一覧と突き合わせ、tarball にないファイルを洗い出した。

\begin{figure}[tbp]
\centering
\begin{tikzpicture}[
  st/.style={draw=cQd,fill=cQf,rounded corners=2pt,text width=3.55cm,align=flush left,font=\footnotesize,
    inner sep=3pt,minimum height=1.45cm,line width=0.5pt},
  num/.style={font=\footnotesize\bfseries,text=cQd!80!black},
]
\node[st] (s1) at (0,0)    {\textbf{1 棚卸し}\\5系統の所在を列挙。展開済みフォルダを元の tarball と照合して自作ファイルを特定};
\node[st] (s2) at (4.25,0) {\textbf{2 退避}\\自作模型・NLOCT の出力・samples を \texttt{models/backup\_20261009} へ \texttt{rsync}};
\node[st] (s3) at (8.5,0)  {\textbf{3 照合}\\\texttt{diff -rq} で全件一致を確認（570ファイル、9.3 MB）};
\node[st] (s4) at (12.75,0){\textbf{4 撤去}\\22項目を \texttt{/usr/bin/trash} でゴミ箱へ。Mathematica には触れない};
\node[st] (s5) at (0,-2.1)    {\textbf{5 取得}\\公式配布元から \texttt{curl} で取得（quarantine 属性が付かない）};
\node[st] (s6) at (4.25,-2.1) {\textbf{6 ビルド}\\FORM、COLLIER、LoopTools、FormCalc の補助ツール};
\node[st] (s7) at (8.5,-2.1)  {\textbf{7 導入}\\Mathematica パッケージ、MadGraph と HEPTools};
\node[st,draw=cGd,fill=cGf] (s8) at (12.75,-2.1){\textbf{8 検証}\\\texttt{tests/run\_all.sh} で10項目};
\foreach \a/\b in {s1/s2,s2/s3,s3/s4,s5/s6,s6/s7,s7/s8} \draw[flow] (\a) -- (\b);
\draw[flow] (s4.south) -- ++(0,-0.18) -| (s5.north);
\end{tikzpicture}
\caption{移行の手順。1〜4 が旧環境の撤去、5〜8 が新しい環境の構築である。自作のファイルは、元と一致することを確かめてから撤去した（手順3）。}
\label{fig:migration}
\end{figure}

\paragraph{見つかった自作ファイル。}突き合わせの結果、FeynArts・FormCalc・MadGraph のフォルダには、ビルドで生じたファイル以外の追加はなかった。MadGraph のフォルダに計算結果は残っていなかった。一方、模型のフォルダには自作のものが多数あった。Library の FeynRules の \texttt{Models} には、AbelianHiggs、B-L-SM、EW、EW+CS、EW+pions、Phi3、QCD、QED、SMnew、SU(2)xU(1)、Yukawa、pseudoYukawa、scalarQED などの模型が356ファイルあった。FeynCalc に同梱された FeynArts の \texttt{Models} には、FeynRules から書き出した模型が15個あった。Documents にあった古いコピーには、これらの旧版があり、QED だけは Library の版と中身が異なっていた。Downloads と Dropbox の \texttt{feynrules-current} には NLOCT の出力（SMQCDrenoL、SM\_NLO）があった。

\paragraph{退避と照合。}これらを元のフォルダ構成のまま \nolinkurl{QFTcalculation/models/backup_20261009/} にコピーした。出どころは同じ場所の README に記録した。コピーのたびに \texttt{diff -rq} で元と比べ、570ファイル（9.3 MB）すべてが一致することを確かめた。今後、自作の模型はインストール先の外のこのフォルダで管理する。FeynRules では \texttt{LoadModel} にフルパスを渡せばよく、模型を FeynRules のフォルダに戻す必要はない。

\paragraph{撤去。}撤去した22項目を表~\ref{tab:removed} に示す。削除には \texttt{rm} ではなく macOS の \texttt{/usr/bin/trash} を使ってゴミ箱へ移した。Finder の「戻す」で元の場所に復元できる。Mathematica の二つの版と、micrOMEGAs 6.0 に同梱された LoopTools 2.12 は、撤去の対象から外した。

\begin{table}[tbp]
\centering\small
\caption{ゴミ箱へ移した22項目。}
\label{tab:removed}
\begin{tabularx}{\linewidth}{>{\raggedright\arraybackslash}p{3.4cm}X>{\raggedleft\arraybackslash}p{1.0cm}}
\toprule
\textbf{場所} & \textbf{撤去したもの} & \textbf{項目数} \\
\midrule
\pth{Library/Wolfram/} \nolinkurl{Applications} & FeynCalc（10.1.0、パッチ版 FeynArts 3.12 同梱）、FeynRules（2.3.49）、LieART（2.1.0）、X（Package-X 2.1.1 の旧配布版） & 4 \\
\pth{Documents} & \texttt{Applications}（FeynCalc・FeynRules・X の旧コピー）、\texttt{feynrules-current} & 2 \\
\pth{Downloads} & FeynArts-3.12、FormCalc-9.10、MG5\_aMC\_v3\_6\_2、feynrules-current、feynrules-samples と、それぞれの tarball・zip & 12 \\
\pth{Dropbox} & GraduateResearch の FeynArts-3.11 と feynrules-current、Examples の FeynRules と feynrules-current & 4 \\
\bottomrule
\end{tabularx}
\end{table}

\section{導入した構成}\label{sec:layout}

導入した環境の全体像を図~\ref{fig:arch} に、各ツールの版と配置を表~\ref{tab:installed} に示す。構成は三つの層に分かれる。一つ目は Mathematica パッケージの層で、Applications に置いた FeynCalc、Package-X、FeynRules からなる。二つ目は C と Fortran の層で、\pth{hep} の下にソースから構築した FORM、LoopTools、COLLIER と、FormCalc 用の FeynArts・FormCalc を置いた。三つ目は MadGraph の層で、conda 環境 \texttt{qft} の Python で動く \MGaMC{} 3.8.1 と、その HEPTools からなる。\pth{hep} の大きさは 549 MB である。

\begin{figure}[tbp]
\centering
\begin{tikzpicture}
\begin{scope}[on background layer]
  \fill[panel] (0,0.35) rectangle (5.4,-5.55);
  \fill[panel] (5.65,0.35) rectangle (11.05,-5.55);
  \fill[panel] (11.3,0.35) rectangle (16.7,-5.55);
\end{scope}
\node[ptitle] at (2.7,0.35) {Mathematica 14.2.1\\Applications};
\node[ptitle] at (8.35,0.35) {\texttt{\textasciitilde/hep}\\（ソースから構築）};
\node[ptitle] at (14.0,0.35) {\MGaMC{} 3.8.1\\conda 環境 \texttt{qft}};
% Mathematica
\node[tM] (FC) at (2.7,-0.95) {FeynCalc 10.2.1};
\node[tM,text width=4.1cm] (FAp) at (2.95,-1.55) {FeynArts 3.12（パッチ版）};
\node[tM,text width=4.1cm] (FH) at (2.95,-2.15) {FeynHelpers 2.0};
\node[tM] (PX) at (2.7,-3.0) {Package-X 2.1.1 + PVReduce 0.1.0};
\node[tN] (ST) at (2.7,-3.8) {入口：FeynArts.m、FormCalc.m};
\node[tM] (FR) at (2.7,-4.75) {FeynRules 2.3.49 + NLOCT\\（局所パッチ 2 つ）};
\draw[cMd,thick] (FC.south west) ++(0.15,0) |- (FAp.west);
\draw[cMd,thick] (FC.south west) ++(0.15,0) |- (FH.west);
% ~/hep
\node[tO] (FORM) at (8.35,-0.95) {FORM 5.0.2 / TFORM};
\node[tO] (LT) at (8.35,-1.75) {LoopTools 2.16\\ライブラリ + MathLink 版};
\node[tO] (COL) at (8.35,-2.6) {COLLIER 1.2.10};
\node[tM] (FAu) at (8.35,-3.35) {FeynArts 3.12（パッチなし）};
\node[tM] (FCa) at (8.35,-4.05) {FormCalc 9.10（FORM 4.2.1 同梱）};
\node[tN] (MCC) at (8.35,-4.8) {代替 \texttt{mcc}（ビルド時のみ）};
% MG5
\node[tO,text width=4.3cm] (MG) at (14.0,-0.95) {\MGaMC{} 3.8.1};
\node[tN,text width=4.3cm] (MOD) at (14.0,-1.75) {模型：sm、loop\_sm、自作の UFO};
\node[tO,text width=4.3cm] (HT) at (14.0,-2.65) {HEPTools：Ninja 1.2.1・\\OneLOop 3.6・COLLIER 1.2.9};
\node[tN,text width=4.3cm] (PY) at (14.0,-3.6) {Python 3.12.15（numpy、scipy、six）};
\draw[flow] (MG) -- (MOD);
% arrows
\draw[flow,cOd] (FH.east) -- (LT.west);
\draw[flow,cMd] (FH.south) -- node[right,font=\scriptsize] {PaXEvaluate} (FH.south |- PX.north);
\draw[flow] (ST.east) -- (FAu.west);
\draw[flow] (ST.east) -- (FCa.west);
\draw[flow,dashed] (MCC.east) -- ++(0.3,0) |- (LT.east);
\draw[flow,cGd] (FR.south) -- ++(0,-0.35) -| node[pos=0.04,above,font=\scriptsize,text=cGd] {WriteUFO} ($(MOD.east)+(0.3,0)$) -- (MOD.east);
% toolchain
\node[draw=black!40,fill=black!6,rounded corners=2pt,text width=16.4cm,align=flush center,font=\footnotesize,inner sep=3pt]
  at (8.35,-6.05) {共通の基盤：Homebrew gcc 15.2（gfortran）・Apple clang 16・macOS SDK 15.2（SDKROOT）・GMP 6.3・MPFR 4.2.2・FLINT 3.6・zstd 1.5.7};
% legend
\begin{scope}[shift={(0.2,-6.85)},every node/.append style={font=\footnotesize}]
  \node[tM,text width=0.5cm,minimum height=0.32cm] (L1) at (0.3,0) {};
  \node[anchor=west] at (L1.east) {Mathematica が必要};
  \node[tO,text width=0.5cm,minimum height=0.32cm] (L2) at (4.3,0) {};
  \node[anchor=west] at (L2.east) {Mathematica は不要};
  \node[tN,text width=0.5cm,minimum height=0.32cm] (L3) at (8.3,0) {};
  \node[anchor=west] at (L3.east) {補助（入口、スクリプト、データ）};
\end{scope}
\end{tikzpicture}
\caption{導入した環境の全体像。左が Mathematica パッケージ（Applications）、中央が \texttt{\textasciitilde/hep} に構築したプログラム、右が MadGraph である。青の矢印は FeynHelpers から LoopTools の MathLink 実行ファイルを呼ぶ経路（\texttt{LToolsEvaluate}）、緑の矢印は FeynRules から MadGraph へ模型を渡す経路である。FeynRules の FeynArts 出力は、パッチ版とパッチなしのどちらの FeynArts でも読める。破線の \texttt{mcc} は LoopTools の構築にだけ使った。}
\label{fig:arch}
\end{figure}

\begin{table}[tbp]
\centering\small
\caption{導入したツールの版、配置、入手元。Applications は \texttt{\textasciitilde/Library/Wolfram/Applications} を表す。}
\label{tab:installed}
\begin{tabularx}{\linewidth}{>{\raggedright\arraybackslash}p{3.3cm}>{\raggedright\arraybackslash}p{1.6cm}>{\raggedright\arraybackslash}X>{\raggedright\arraybackslash}p{3.3cm}}
\toprule
\textbf{ツール} & \textbf{版} & \textbf{配置} & \textbf{入手元} \\
\midrule
FORM / TFORM & 5.0.2 & \texttt{\textasciitilde/hep/bin} & GitHub（form-dev）~\cite{form5,url:form} \\
LoopTools & 2.16 & \texttt{\textasciitilde/hep/\{lib,include,bin\}} & feynarts.de~\cite{fcalc,url:feynarts} \\
COLLIER & 1.2.10 & \texttt{\textasciitilde/hep/lib} & HepForge~\cite{collier,url:collier} \\
FeynCalc & 10.2.1 & Applications/FeynCalc & 公式インストーラ~\cite{fc102,fc10} \\
FeynArts（パッチ版） & 3.12 & FeynCalc/FeynArts & FeynCalc のミラー~\cite{fa3} \\
FeynHelpers & 2.0.0 & FeynCalc/AddOns/FeynHelpers & 公式インストーラ~\cite{fc102,fh1} \\
Package-X（ソース版） & 2.1.1 & Applications/X & GitLab（mule-tools）~\cite{px1,px2,url:px} \\
PVReduce & 0.1.0 & Applications/X/PVReduce.m & 同上 \\
FeynArts（パッチなし） & 3.12 & \texttt{\textasciitilde/hep/mathematica} & feynarts.de~\cite{fa3,url:feynarts} \\
FormCalc & 9.10 & \texttt{\textasciitilde/hep/mathematica} & feynarts.de~\cite{fcalc,url:feynarts} \\
FeynRules + NLOCT & 2.3.49 & Applications/FeynRules & 公式サイト~\cite{fr2,nloct,url:fr} \\
\MGaMC{} & 3.8.1 & \texttt{\textasciitilde/hep/MG5\_aMC\_v3\_8\_1} & Launchpad~\cite{mg5,url:mg5} \\
HEPTools（Ninja、OneLOop、COLLIER） & 1.2.1、3.6、1.2.9 & MG5 の \texttt{HEPTools} & MadGraph の \texttt{install}~\cite{ninja,ninjamg,oneloop,collier} \\
\bottomrule
\end{tabularx}
\end{table}

\paragraph{FeynArts を二つ置いた理由。}FeynCalc は、自分の記号とぶつかる FeynArts の記号を改名するパッチを当てて、FeynArts を読み込む。一方、FormCalc はパッチのない FeynArts を前提にしている。そこで、パッチ版は FeynCalc のフォルダの中に、パッチなし版は \pth{hep/mathematica} に置いた。FeynCalc はパッチ版を自分のフォルダから明示的なパスで読むので、二つは衝突しない。パッチなし版と FormCalc を \texttt{<<FeynArts\bq} や \texttt{<<FormCalc\bq} で読めるように、Applications には数行の入口ファイル（\ref{app:files}）を置いた。同じカーネルで両方の FeynArts を読み込んではならない。

\paragraph{Package-X の版。}Package-X は開発元が更新をやめており、McMule グループが CC-BY-4.0 で再公開している~\cite{url:px}。旧配布版は Mathematica の \texttt{Encode} で難読化されていたが、今回は同じ 2.1.1 の読めるソースを入れた。PVReduce のコンテキストは \texttt{X\bq PVReduce\bq} なので、ファイルは \texttt{X/PVReduce.m} に置いた。COLLIER を Package-X から呼ぶ拡張 CollierLink は導入できなかった。実行時に必要なライブラリが Intel Mac 用しか配布されていないためである。

\paragraph{COLLIER が二つある理由。}単体の COLLIER（1.2.10）と、MadGraph が HEPTools として自前で入れた COLLIER（1.2.9）は、別々のライブラリである。前者は Fortran や C から直接使い、後者は MadLoop だけが使う。互いに干渉しない。

\section{この環境に固有の問題と対処}\label{sec:issues}

導入の途中で見つかった問題と、その対処を表~\ref{tab:issues} にまとめる。原因が確定していないものは「推論」と書いた。以下、主なものを順に説明する。

\begin{table}[p]
\centering\footnotesize
\caption{この Mac に固有の問題と対処。}
\label{tab:issues}
\begin{tabularx}{\linewidth}{>{\raggedright\arraybackslash}p{2.4cm}>{\raggedright\arraybackslash}X>{\raggedright\arraybackslash}X>{\raggedright\arraybackslash}X}
\toprule
\textbf{対象} & \textbf{症状} & \textbf{原因} & \textbf{対処} \\
\midrule
gfortran のリンク & \texttt{ld: library 'System' not found} & gfortran 15.2 が macOS SDK を見つけない & \texttt{\textasciitilde/.zshrc} で SDKROOT を設定 \\
MathLink の \texttt{mcc} & パスが空白で分断され、i386 向けにリンクされる & 2012年版のスクリプトが arm64 とパスの空白を想定していない & 代替の \texttt{\textasciitilde/hep/bin/mcc} \\
FeynRules & 読み込み時に \texttt{SetDelayed::write}（\texttt{MatrixSymbol}） & 14.1 で入った \texttt{System\bq MatrixSymbol} が同名の記号を隠す & FeynRules.m に局所パッチ（元は \texttt{.orig}） \\
NLOCT の繰り込み & \texttt{OnShell\allowbreak Renormalization} が、場の混合のない模型で止まる（\texttt{GetFC::Head}） & 混合の項を除く \texttt{DeleteCases} が、項が1つのとき積の因子を消す & \texttt{Core/Loop.m} に局所パッチ（元は \texttt{.orig}） \\
NLOCT の $\delta Z$ & 場の繰り込みが決まらず、記号のまま出力に残る & flavor を展開しないと \texttt{FR\$Classes\allowbreak Translation} が空で、自分の場の $\delta Z$ を0と置く & \texttt{WriteFeynArtsOutput[..., FlavorExpand -> True]} \\
NLOCT の出力 & R2 がないと \texttt{R2\$vertlist = \{Last[\{\}]\}}。\texttt{MSbar -> True} では2点の対抗項の $-i$ が抜ける & NLOCT の不具合 & 読み込み後に消す。$\overline{\mathrm{MS}}$ の対抗項は on-shell の極から作る \\
MadGraph の1ループ & BSM の結合の次数でループを生成しない。\texttt{gauge = [0]} の模型も拒む & ループ用インターフェースが QCD と QED しか検査しない & 作業用の UFO で次数を QCD と改名。\texttt{FeynmanGauge = True} \\
MadLoop の検査プログラム & 点ごとに $\mu_R=\sqrt s$ に置き直す & \texttt{check\_sa.f} の仕様 & 生成されたプログラムを書き換え、ファイルから読む \\
MadGraph の木の積分 & $\phi\phi\to\phi\phi\phi$ の断面積が、$\sqrt s=10m$ で回ごとに表示の誤差より大きくばらつき、自前の積分より $0.8\%$ 低い & 不明。行列要素は一致するので、位相空間の積分（推論） & 重点サンプリングの自前の積分と比べる。複数の乱数の種で回す \\
FeynHelpers & \texttt{LToolsLoadLibrary\allowbreak ::nobinary} & LoopTools を既定の場所でしか探さない & シンボリックリンクを置く \\
ブラウザで取得した実行ファイル & 起動した直後にファイルが消える & quarantine 属性の付いた未署名のバイナリ（推論） & \texttt{curl} で取得し直す \\
CollierLink & 導入できない & 実行時ライブラリが Intel Mac 用のみで、ソースが非公開 & FeynHelpers の \texttt{LToolsEvaluate} で代用 \\
MadGraph の \texttt{install} & HEPTools が作業フォルダに入る & 相対パスで導入される & MadGraph のフォルダで実行 \\
Ninja & 4倍精度なしでビルドされる & HEPTools の構成 & MadLoop は不安定な点で COLLIER・CutTools に切り替える \\
LoopTools の $C_0$ & 閾値より下で外運動量がゼロでない $C_0$ が誤る（$C_0(1,1,1)$ が $-0.671$ でなく $6.87+34.2i$）。Fortran から呼ぶと止まる & 不明。gfortran 14 と 15、最適化の有無によらない & $C_0$ には Package-X の \texttt{ScalarC0} と COLLIER を使う。$A_0$、$B_0$ は正しい \\
LoopTools の $C_{00}$ & 閾値より下で外運動量がゼロでない $C_{00}$ が誤る。外運動量がすべてゼロでは \texttt{NaN} を返す（Fortran から呼んでも同じ） & $C_{00}$ は $C_0$ を含む（Passarino--Veltman の簡約）ので、$C_0$ の誤りが及ぶ（推論） & $C_{00}$ には Package-X と COLLIER を使う。COLLIER の $C_{00}$ は外運動量がゼロでも正しい \\
COLLIER の $C_0$ & 外運動量がすべて厳密にゼロのとき $10^{50}$ を返す & 仕様（``case not implemented''） & Package-X か LoopTools を使う \\
FeynHelpers の $C_0$ 展開 & \texttt{PaXC0Expand -> True} の式が閾値より下で \texttt{Undefined} & 適用範囲付きの式（\texttt{ConditionalExpression}） & Package-X の \texttt{X\bq ScalarC0} で数値評価する \\
$D_0$（外線と内線の質量がすべて等しい場合） & LoopTools は $t=0$ で \texttt{NaN}。COLLIER は $t=0$ で $10^{50}$ を返し、$s,t,u$ がすべて $4m^2$ より小さい点では停止する。Package-X の \texttt{X\bq ScalarD0} はその点で \texttt{Indeterminate} & 退化した運動学を各実装が扱わない（FF：``on threshold''、COLLIER：``lambda<0''） & どのライブラリも値を返す点で比べる。物理的な散乱の点（$s>4m^2$、$t,u<0$、$t\ne0$）では3つとも一致する \\
COLLIER の $D_{00}$（$s\gg m^2$） & COLI は虚部が徐々にずれ（$\sqrt s=10^4m$ で $44\%$）、DD は $\sqrt s\gtrsim1.3\times10^3m$ で全体が誤る & 不明 & LoopTools か、Package-X の解析式（\texttt{PaXD0Expand -> True}）の高精度の値で確かめる \\
FeynHelpers の \texttt{LToolsEvaluate} & 倍精度の係数を持つ $B_0$ に対して、テンソル積分を無視したという警告（\texttt{::tens}）を出す & 誤判定（値は正しい） & 運動学を有理数のまま渡す \\
miniforge & base 環境で numpy を読み込めない。\texttt{mamba} は \texttt{ImportError} & \texttt{libgfortran} の rpath の重複。古い mamba と新しい conda（推論） & 環境 \texttt{qft} と \texttt{conda install} を使う \\
\bottomrule
\end{tabularx}
\end{table}

\paragraph{gfortran と macOS SDK。}LoopTools の configure は、Fortran と C のリンクを確かめる段階で止まった。同じ小さなプログラムを手でリンクすると、\texttt{ld: library 'System' not found} が出た。Homebrew の gfortran 15.2 が、macOS 26.2 の Command Line Tools にある SDK（15.2）を見つけられないためである。環境変数 \texttt{SDKROOT=\$(xcrun --show-sdk-path)} を与えるとリンクが通り、テストプログラムは期待どおりの終了コード 95 を返した。MadGraph はプロセスを生成するたびに gfortran でコンパイルするので、この設定は今後も常に必要になる。そこで \pth{.zshrc} の末尾に専用のブロックを足した（\ref{app:files}）。ただし、Dock から起動したアプリは \pth{.zshrc} を読まない。Mathematica のノートブックから Fortran をコンパイルする場合は、ターミナルで行う必要がある。

\paragraph{MathLink の \texttt{mcc}。}LoopTools の Mathematica インターフェースは、MathLink の実行ファイルとして作られる。そのためのコンパイラ・スクリプト \texttt{mcc} は、Mathematica 14.2.1 に付属するものでも2012年版のままだった。このスクリプトは \texttt{Wolfram 2.app} の空白でパスを分断し、arm64 の Mac でも i386 向けにリンクしようとする。64ビットを指定する \texttt{-b64} も「この環境では生成できない」として拒否された。一方、部品である \texttt{mprep} と \texttt{libMLi4.a}（x86\_64 と arm64 のユニバーサル版）はそろっていた。そこで、\texttt{.tm} ファイルを \texttt{mprep} で C に変換し、arm64 の \texttt{libMLi4.a} を静的にリンクするだけの代替スクリプトを書いた（\ref{app:files}）。これで作った試験用の MathLink プログラムは、Mathematica から呼ぶと \texttt{Test[41] = 42} を返した。LoopTools はこのスクリプトで構築した。

\paragraph{FeynRules と \texttt{MatrixSymbol}。}FeynRules 2.3.49（2021年）は今も最新版である。Mathematica 14.1 は \texttt{System\bq MatrixSymbol} を組み込み記号として加えた（\texttt{WolframLanguageData} の導入版は 14.1）。FeynRules は同じ名前の記号を自分のコンテキストで定義し、模型ファイルの混合行列の宣言でも使う。組み込み記号が先に見つかると、FeynRules の定義は保護された記号への代入とみなされ、\texttt{SetDelayed::write} で失敗する。問題を難しくしていたのは、FeynRules が Feynman 則の計算を並列化するために8個のサブカーネルを立ち上げ、それぞれで FeynRules を読み直すことである（図~\ref{fig:matrixsymbol}）。主カーネルでだけ組み込み記号を外す試験では、主カーネルのメッセージは5件から0件に減ったが、サブカーネルから同じメッセージが24件届いた。そこで FeynRules.m の2か所、すなわち主カーネルでの読み込みの直前と、サブカーネルへの \texttt{ParallelEvaluate} の中で、組み込み記号を外すようにした。コードを文字列で持たせ、各カーネルで \texttt{ToExpression} する形にしたのは、ファイルを読む時点で記号が解決されてしまうのを避けるためである。パッチの後は、標準模型の読み込み、UFO 出力、FeynArts 出力を通してメッセージは0件になった。

\begin{figure}[tbp]
\centering
\begin{tikzpicture}[
  k/.style={draw=cPd,fill=white,rounded corners=2pt,text width=6.6cm,align=flush center,font=\footnotesize,inner sep=3pt,minimum height=0.95cm},
  bad/.style={draw=cRd,fill=cRf,rounded corners=2pt,text width=6.6cm,align=flush center,font=\footnotesize,inner sep=3pt},
  good/.style={draw=cGd,fill=cGf,rounded corners=2pt,text width=6.6cm,align=flush center,font=\footnotesize,inner sep=3pt},
  hd/.style={font=\small\bfseries},
]
\node[hd] at (0,0.95) {パッチ前};
\node[hd] at (8.6,0.95) {パッチ後};
% before
\node[k] (m0) at (0,-0.3) {主カーネル：\texttt{Get[FeynRulesPackage.m]}\\\texttt{MatrixSymbol} $\to$ \texttt{System\bq MatrixSymbol}（14.1〜、保護）};
\node[k,fill=black!4] at (0.12,-1.68) {};
\node[k,fill=black!2] at (0.06,-1.62) {};
\node[k] (s0) at (0,-1.56) {並列サブカーネル ×8：\texttt{ParallelEvaluate[Get[...]]}\\同じく組み込み記号に解決される};
\node[bad] (r0) at (0,-3.0) {FeynRules の \texttt{MatrixSymbol} の定義が\\各カーネルで \texttt{SetDelayed::write} により失敗};
\draw[flow,cRd] (m0) -- (s0);
\draw[flow,cRd] (s0) -- (r0);
% after
\node[k,draw=cMd,fill=cMf] (p) at (8.6,0.0) {FeynRules.m：\texttt{FR\$MatrixSymbolFix}（文字列）};
\node[k] (m1) at (8.6,-1.0) {主カーネル：\texttt{ToExpression[FR\$MatrixSymbolFix]}\\$\to$ \texttt{Get[FeynRulesPackage.m]}};
\node[k] (s1) at (8.6,-2.1) {サブカーネル ×8：\texttt{ParallelEvaluate} 内で同じ処理\\$\to$ \texttt{Get[FeynRulesPackage.m]}};
\node[good] (r1) at (8.6,-3.2) {\texttt{MatrixSymbol} $\to$ \texttt{FeynRules\bq MatrixSymbol}\\標準模型の読み込みと出力でメッセージ 0 件};
\draw[flow] (p) -- (m1);
\draw[flow] (m1) -- (s1);
\draw[flow,cGd] (s1) -- (r1);
\end{tikzpicture}
\caption{FeynRules と \texttt{MatrixSymbol} の衝突と、局所パッチの働き。FeynRules はサブカーネルでも自分を読み込み直すので、組み込み記号は主カーネルとサブカーネルの両方で外す必要がある。主カーネルだけで外した試験では、サブカーネルから24件のメッセージが届いた。}
\label{fig:matrixsymbol}
\end{figure}

FeynMaster のマニュアルは、FeynRules の3つのファイルで \texttt{MatrixSymbol} を \texttt{FRMatrixSymbol} に改名する対処を、FeynRules の開発者の提案として紹介している~\cite{url:feynmaster}。今回はこの方法を採らなかった。\texttt{MatrixSymbol} は模型ファイルの混合行列の宣言で利用者が書く名前なので、改名するとそれを使う模型が動かなくなるからである。退避した自作の模型には \texttt{MatrixSymbol} を使うものはなかった。また、ほかにも同じ種類の衝突がないかを確かめた。FeynRules のソースに現れる識別子のうち組み込み記号と同じ名前のもの465個について、導入された版を \texttt{WolframLanguageData} で調べたところ、13.0以降に導入されたものは \texttt{MatrixSymbol}（14.1）だけだった。なお、FeynRules 2.3.49 を Mathematica 14.0.0 で使って UFO を生成した例が、2024年の論文の付属データにある~\cite{allanach}。

\paragraph{FeynHelpers が LoopTools を探す場所。}FeynHelpers の \texttt{LToolsEvaluate} は、LoopTools の MathLink 実行ファイルを \nolinkurl{FeynHelpers/ExternalTools/LoopTools/LoopTools} にしか探さない。ここに \pth{hep/bin/LoopTools} へのシンボリックリンクを置いた。ビルドしたものは一つだけなので、LoopTools を更新してもリンクを張り直す必要はない。

\paragraph{ブラウザで取得した実行ファイル。}旧環境で FormCalc に付属する FORM のバージョンを確かめるために起動したところ、二か所にあった \texttt{form} が直後に消えた。ゴミ箱にはなく、読めた範囲のシステムログにも記録はなかった。どちらのファイルにも quarantine 属性が付いていたので、macOS のセキュリティ機能が未署名のバイナリを削除したと考えられる（推論）。今回は \texttt{curl} で取得し直した。そのため、ファイルに付くのは \texttt{com.apple.provenance} だけで、同じ FormCalc 付属の FORM 4.2.1 は問題なく起動した。

\paragraph{CollierLink、HEPTools、Ninja。}CollierLink は実行時に \texttt{libcollierX.dylib} を呼ぶが、このライブラリは Intel Mac 用しか配布されておらず、ソースも公開されていない。Package-X の式の数値評価は、FeynHelpers 経由の LoopTools で代用できる。MadGraph の \texttt{install ninja}・\texttt{install collier} は、実行したときの作業フォルダに HEPTools を作る。MadGraph のフォルダで実行し直して、正しい場所に入れた。Ninja は4倍精度なしでビルドされたが、MadLoop は数値が不安定な位相空間点で COLLIER や CutTools に切り替えるので、実用上の支障はない。

\paragraph{LoopTools の $C_0$ と関連する制限。}これは環境を組んだ後、$\phi^3$ 模型の $d=4$ の1ループの検算（\nolinkurl{models/phi3/ResearchNote/note.pdf}）で見つかった。LoopTools 2.16 は、閾値より下で外運動量がゼロでない $C_0$ を誤って返す。例えば内線の質量がすべて1の $C_0(1,1,1)$ と $C_0(1,1,3.5)$ は、Feynman パラメータの2重積分では $-0.671253105748$ と $-1.101657275584$ の実数になる。ところが MathLink 経由の LoopTools は $6.872+34.19i$ と $10.49+6.87i$ を返す。LoopTools の比較モード（\texttt{SetVersionKey}）では、内部の二つの実装の一方は正しい値を出すが、どの版を選んでも誤った値が返ってくる。Fortran から直接呼ぶと計算が終わらず、$(0.5,0.7,2)$ のような点でも同じだった。最適化なし（\texttt{--debug}、\texttt{-O0}）でビルドしても、Homebrew に残っていた gfortran 14.2 でビルドしても変わらない。一方、$A_0$、$B_0$、外運動量がゼロの $C_0$、閾値より上の $C_0$ は正しく、$B_0$ は15点で Package-X・COLLIER・閉じた式と相対 $4.4\times10^{-15}$ 以内で一致した。原因は突き止めていないので、当面 LoopTools の $C_0$ は使わない。COLLIER の $C_0$ は、外運動量がすべて厳密にゼロの場合を「case not implemented」として扱わず、番兵値 $10^{50}$ を返す（運動量をわずかにずらせば正しい極限に近づく）。FeynHelpers で Package-X の $C_0$ を展開した式（\texttt{PaXC0Expand -> True}）は、閾値より下で \texttt{Undefined} になる適用範囲付きの式である。Package-X 本体の \texttt{X\bq ScalarC0} は全点で正しく評価できた。

続く $d=6$ の1ループの検算では、$D+2$ 次元のスカラー積分が $D$ 次元のテンソル積分の $g^{\mu\nu}$ の係数の $-2\pi$ 倍であること（FeynCalc の規格化で $\mathrm{PaVe}^{(D+2)}=-2\pi\,\mathrm{PaVe}_{00}^{(D)}$）を使い、$D=4-2\epsilon$ 用のツールの $A_{00}$、$B_{00}$、$C_{00}$ から6次元の積分を求めた。そこで、LoopTools の $C_{00}$ も閾値より下で外運動量がゼロでない点では誤ることが分かった。例えば $C_{00}(1,1,1)$ から求めた6次元の有限部分は、正しい値 $-0.1456$ に対して $-5.17-22.8i$ になる。外運動量がすべてゼロの点では、$C_0$ は正しい $-1/2$ を返すのに、$C_{00}$ は Fortran から直接呼んでも \texttt{NaN} を返す。Passarino--Veltman の簡約では $C_{00}$ は $C_0$ を含むので、$C_0$ の誤りが及んだと考えられる。LoopTools の $A_{00}$ と $B_{00}$ は正しく、COLLIER の $C_{00}$ は外運動量がすべてゼロの点も含めて正しかった。

さらに、1ループの $\phi\phi\to\phi\phi$ 振幅の検算で、4点関数 $D_0$ を扱える範囲がライブラリごとに違うことが分かった。外線の $p_i^2$ と内線の質量がすべて $m$ の運動学で、LoopTools は前方 $t=0$ で値を返さない（FF の実装が「閾値上の退化した場合」として扱えず、\texttt{NaN} になる）。COLLIER も $t=0$ では番兵値 $10^{50}$ を返す。$s,t,u$ がすべて $4m^2$ より小さい非物理的な点（例えば $s=2m^2$、$t=-m^2$）では、COLLIER は致命的エラーで停止し（\texttt{D0comb\_coli: case 1, lambda<0}）、Package-X の数値関数 \texttt{X\bq ScalarD0} は \texttt{Indeterminate} を返す。\texttt{X\bq ScalarD0} は厳密な数を引数にすると評価されずに残り、それに \texttt{N[..., 20]} を掛けると計算が終わらない。物理的な散乱の点（$s>4m^2$、$t,u<0$、$t\ne0$）では、3つのライブラリの $D_0$ は全桁で一致した。この振幅に現れる $C_0(m^2,m^2,x)$（$x<0$ または $x>4m^2$）は、LoopTools でも正しかった。また、FeynHelpers の \texttt{LToolsEvaluate} は、倍精度の数を係数に持つ $B_0$ を含む式に対して「テンソル積分を無視した」という警告（\texttt{::tens}）を出すが、値は正しい。運動学を有理数のまま渡せば警告は出ない。

\paragraph{NLOCT と MadLoop。}$\phi^3$ 模型の1ループの $\phi\phi\to\phi\phi$ を MadLoop と比べるため（\nolinkurl{models/phi3/ResearchNote/note.pdf} の手順 (6)）、FeynRules に付属する NLOCT 1.02 で対抗項を含む UFO を作った。NLOCT の関数 \texttt{OnShellRenormalization} は、場の混合を扱う段階で、2点の対抗項の項のうち同じ場の組のものを \texttt{DeleteCases} で除く。混合のない模型ではこの段階に渡る式が1つの積になり、\texttt{DeleteCases} がその因子をすべて消して \texttt{FeynmanRules[1]} を呼び、\texttt{GetFC::Head} で止まる。FeynRules の \texttt{Core/Loop.m} に局所パッチを当て、項のリストに対して \texttt{DeleteCases} を適用し、混合の項が残らないときはこの段階を飛ばすようにした（\ref{app:files}）。元のファイルは \texttt{Loop.m.orig} として残した。次に、FeynRules は flavor を展開しないと FeynArts の模型の対応表 \texttt{FR\$ClassesTranslation} を空にする。すると NLOCT は FeynArts の粒子 \texttt{S[1]} を模型の場と結びつけられず、その場自身の $\delta Z$ を「関係のない場の $\delta Z$」として0と置き、$\delta Z$ が決まらないまま出力に残る。FeynArts の模型を \texttt{FlavorExpand -> True} で書き出すと対応表が書かれ、この問題は消える。NLOCT はさらに、R2 の頂点がないと不正な \texttt{R2\$vertlist = \{Last[\{\}]\}} を書き、\texttt{MSbar -> True} では2点の対抗項から因子 $-i$ を落とす。MadGraph 3.8.1 の1ループ用インターフェースは、ループ補正の種類を QCD か QED としてしか検査しない（\texttt{validate\_model} の既定値が QCD）。そのため、BSM の結合の次数（ここでは \texttt{PHI}）を摂動に指定すると「この模型はループを扱えない」として止まる。また、UFO が \texttt{gauge = [0]}（ユニタリーゲージのみ）と宣言していると、QCD 以外のループを生成しない。作業用の UFO のコピーで次数の名前を \texttt{QCD} に変え、\texttt{FeynmanGauge = True} として \texttt{gauge = [0, 1]} と書かせて回避した。MadLoop の単体の検査プログラム \texttt{check\_sa.f} は、位相空間の点ごとに $\mu_R=\sqrt s$ と置き直すので、比較では生成されたプログラムを書き換えて $\mu_R$ をファイルから読ませた。これらの回避策のもとで、NLOCT の対抗項は手計算の on-shell の処方と記号として一致し、MadLoop の値は手計算の式と相対 $3\times10^{-13}$ 以内で一致した。

\paragraph{MadGraph の木の $2\to3$ の積分。}同じ $\phi^3$ 模型の全断面積の計算（\nolinkurl{models/phi3/ResearchNote/note.pdf} の手順 (7)）で、$\phi\phi\to\phi\phi\phi$ の木の断面積を MadGraph（固定エネルギーのビーム、カットなし）と自前のモンテカルロ積分で比べた。MadGraph の行列要素は、自前の式と同じ位相空間の点で相対 $10^{-15}$ まで一致した。断面積は $\sqrt s=4m$ と $5m$ では統計誤差の0.5倍以内で一致した。ところが $\sqrt s=10m$ では、乱数の種、事象数（1万と5万）、積分の設定（\texttt{sde\_strategy} の1と2、\texttt{hard\_survey 2}）を変えた5回の結果が $9.165$〜$9.260$~pb とばらついた。このばらつき（標準偏差 $0.036$~pb）は表示された誤差（約 $0.01$~pb）より大きく、平均 $9.220$~pb は、位相空間の変数の取り方が異なる自前の二つの積分（$9.295\pm0.005$ と $9.305\pm0.008$~pb、どちらも VEGAS 型の重点サンプリング）より $0.8\%$ 低い。高いエネルギーでは、二つの粒子がほぼビームの向きに出る狭い領域が断面積の大部分を担う。自前の積分でも、一様なサンプリングでは同じように誤差を小さく見積もった。MadGraph のずれも、この領域の扱いから来ていると考えられるが、原因は突き止めていない。MadGraph の断面積を検算に使うときは、複数の乱数の種で回してばらつきを確かめ、可能なら独立な積分と比べる。

\paragraph{高エネルギーでの COLLIER の箱図。}$\phi^3$ 模型の $d=6$ の漸近自由を散乱で見る計算（\nolinkurl{models/phi3/ResearchNote/note.pdf} の手順 (11)）で、90°散乱の1ループの振幅を $\sqrt s=10^4m$ まで求めた。外線と内線の質量はすべて $m$ で、振幅には $D_{00}$、$C_{00}$、$B_{00}$ が現れる。基準には、FeynHelpers で Package-X の $C_0$ と $D_0$ を対数と二重対数の式に展開し（\texttt{PaXC0Expand -> True}、\texttt{PaXD0Expand -> True}）、厳密な引数で30桁で評価した値を使った。$C_0(m^2,m^2,x)$ の展開式は、$x<0$、$x=m^2$、$x>4m^2$ のいずれでも正しい値を与えた（表~\ref{tab:issues} の \texttt{Undefined} になる場合とは異なる）。LoopTools はこの基準と、$\sqrt s=10^4m$ まで相対 $2.4\times10^{-13}$ 以内で一致した。COLLIER は、$s$ チャンネルを含む箱図 $D_{00}(s,t)$ で精度を失った。既定の COLI（モード1）では、実部は $10^{-14}$ で正しいが、虚部のずれが $\sqrt s\approx130m$ から $10^{-8}$ を超え、$830m$ で $7\times10^{-5}$、$10^4m$ で $44\%$ になった。DD（モード2）は $\sqrt s\approx830m$ までは $10^{-8}$ 程度で正しいが、$1.26\times10^3m$ 以上では実部も虚部も全く誤った値を返した。警告やエラーは出なかった。二つの方式を比べれば食い違いには気づけるが、どちらが正しいかは決まらない。高エネルギーや、質量に比べて不変量が大きい計算では、LoopTools か Package-X の高精度の値で確かめる必要がある。COLLIER の精度の目安を返す関数（\texttt{GetAccFlag\_cll} など）は試していない。

同じ一連の検算で、三つの小さな注意点も分かった。第一に、FeynCalc の PaVe 関数の $\epsilon^0$ の係数は $\overline{\mathrm{MS}}$ の値から極の係数の $-(\gamma_E+\ln\pi)$ 倍だけずれる（図~\ref{fig:b0}）。紫外で有限な $D_{00}$ では差が出ないが、$B_{00}$ と $C_{00}$ では実部がずれる。第二に、Package-X の結果を運動量を記号のまま簡約してから $p^2=0$ を代入すると、$0/0$ になって値が出ない。運動量の値を代入してから \texttt{PaXEvaluate} を呼ぶ。第三に、LoopTools の $C_{00}$ は閾値より下で値が誤るが、その $\mu$ 依存は正しかった。$\mu$ 依存は紫外の対数から来て、それは $B_0$ にしか含まれないからである。なお、\texttt{wolframscript} の終了時に MathLink の例外のメッセージ（\texttt{MLException}）が出ることがあるが、結果のファイルは書き終わっており、値に影響はない。

\paragraph{検証で踏んだ二つの落とし穴。}一つ目は、FeynRules の試験で Feynman 則の計算結果を変数 \texttt{v} に代入したことである。\texttt{v} は標準模型ファイルでヒッグス場の真空期待値を表すので、続く UFO 出力からフェルミオンとゲージ粒子の頂点がすべて消え、評価されないコードを粒子名に含む頂点が3000個以上できた。一時は FeynRules と Mathematica 14 の非互換だと判断しかけたが、代入をやめると出力は正常に戻った。模型の記号に値を代入してはならない。二つ目は、試験スクリプトに \texttt{formcalc.wl} という名前を付けたことである。macOS のファイル名は大文字と小文字を区別しないので、そのフォルダで \texttt{Get["FormCalc\bq"]} を実行すると、Mathematica は本物の FormCalc ではなくこのスクリプト（\texttt{FormCalc.wl} と同一視される）を読み、自分自身を呼んで再帰の上限に達した。試験スクリプトにはすべて \texttt{test\_} を付けた。

\section{検証}\label{sec:tests}

各ツールが正しく動くことを、表~\ref{tab:tests} の10項目で確かめた。項目ごとに、結果を独立に知る手段を用意した。手段は三種類ある。一つ目は既知の厳密値で、$B_0$ と $C_0$ がこれに当たる。二つ目は独立な別実装で、同じ $B_0$ を四つの経路で計算した。三つ目は独立な解析計算で、$e^+e^-\to\mu^+\mu^-$ の断面積がこれに当たる。厳密値と解析式の導出は\ref{app:derivations} に、試験スクリプトは \nolinkurl{environment/tests/} にある。

\begin{table}[tbp]
\centering\small
\caption{検証の結果。「期待値」は独立に求めた値、「結果」は導入した環境での値である。MadGraph の誤差はモンテカルロ積分の統計誤差。}
\label{tab:tests}
\begin{tabularx}{\linewidth}{>{\raggedright\arraybackslash}p{0.35cm}>{\raggedright\arraybackslash}p{2.7cm}>{\raggedright\arraybackslash}X>{\raggedright\arraybackslash}X>{\raggedright\arraybackslash}p{1.9cm}}
\toprule
& \textbf{対象} & \textbf{期待値} & \textbf{結果} & \textbf{判定} \\
\midrule
1 & FORM 5.0.2 & $(x+y)^4$ の二項展開 & $y^4+4xy^3+6x^2y^2+4x^3y+x^4$ & 一致 \\
2 & COLLIER 1.2.10 & $B_0(2;1,1)=2-\pi/2=0.429203673205103$ & $0.429203673205$ & 出力12桁で一致 \\
3 & LoopTools 2.16（MathLink） & $B_0(2;1,1)=0.42920367320510338$、$C_0(0,0,0;1,1,1)=-1/2$ & $0.42920367320510344$、$-0.5$ & 倍精度で一致 \\
4 & FeynHelpers 2 & LoopTools 経由と Package-X 経由の $B_0$ が一致 & ともに $1/\epsilon-1.2927418775458297$ & 一致 \\
5 & Package-X + PVReduce & $B_0=2+1/\epsilon+\mathrm{DiscB}+\ln(\mu^2/m^2)$；$B_1=-B_0/2$（等質量） & 同じ式 & 一致 \\
6 & FeynCalc + FeynArts（パッチ版） & 光子の自己エネルギー：9種のフェルミオンごとに $A_0$ と $B_0$、横波の構造 & PaVe 関数18個、$p^2g^{\mu\nu}-p^\mu p^\nu$ の構造あり & 期待どおり \\
7 & FeynArts（パッチなし）+ FormCalc & $e^+e^-\to\mu^+\mu^-$（木）で \texttt{CalcFeynAmp} が完了 & 振幅1個、略記8個 & 完了 \\
8 & FeynRules 2.3.49 & 標準模型：メッセージなし、MadGraph の \texttt{sm} の頂点をすべて含む & メッセージ0件、163頂点が \texttt{sm} の153頂点をすべて含む & 一致 \\
9 & MadGraph（木） & $e^+e^-\to\mu^+\mu^-$、$\sqrt{s}=1$~TeV、$|\eta_\ell|<2.5$：解析式 0.10426 pb & $0.10437\pm0.00018$ pb & $0.6\sigma$ で一致 \\
10 & MadGraph（ループ誘起） & $gg\to H$、13 TeV、LO：独立な検算なし & $15.76\pm0.05$ pb & 実行を確認 \\
\bottomrule
\end{tabularx}
\end{table}

\paragraph{$B_0$ の相互検算。}最も重要な検算は、同じ $B_0(2;1,1)$ を独立な経路で計算させることである（図~\ref{fig:b0}）。厳密値は $2-\pi/2=0.42920367320510338\ldots$ である（\ref{app:derivations}）。LoopTools はこれを倍精度の全桁で再現し、COLLIER は出力した12桁で再現した。Package-X の解析式 $2+1/\epsilon+\mathrm{DiscB}(s;m,m)+\ln(\mu^2/m^2)$ は、$s=2$、$m=\mu=1$ で $\mathrm{DiscB}=-\pi/2$ となり、同じ有限部分を与える。FeynHelpers を通すと、LoopTools 経由でも Package-X 経由でも $1/\epsilon-1.2927418775458297$ が返り、両者は倍精度の全桁で一致した。この値は $2-\pi/2-\gamma_E-\ln\pi$ に等しく、厳密値とは $-(\gamma_E+\ln\pi)$ だけずれる。このずれは、FeynCalc が紫外発散を $1/\epsilon-\gamma_E-\ln\pi$ の形で書く規格化をとるためと解釈できる（推論）。いずれにせよ、二つの経路が同じ規格化で同じ値を返すことが確かめられた。

\begin{figure}[tbp]
\centering
\begin{tikzpicture}[
  n/.style={draw=cOd,fill=cOf,rounded corners=2pt,text width=4.9cm,align=flush center,font=\footnotesize,inner sep=3pt},
  nm/.style={draw=cMd,fill=cMf,rounded corners=2pt,text width=4.9cm,align=flush center,font=\footnotesize,inner sep=3pt},
  ex/.style={draw=cGd,fill=cGf,rounded corners=3pt,text width=4.2cm,align=flush center,font=\footnotesize,inner sep=4pt,line width=0.8pt},
  e/.style={thick,black!60},
  lab/.style={font=\scriptsize,fill=white,inner sep=1pt,align=flush center},
]
\node[ex] (E) at (0,0) {\textbf{厳密値}\\$B_0(2;1,1)=2-\pi/2$\\$=0.42920367320510338$};
\node[n] (LT) at (-5.4,1.7) {LoopTools 2.16（MathLink）\\$0.42920367320510344$};
\node[n] (CO) at (5.4,1.7) {COLLIER 1.2.10（Fortran）\\$0.429203673205$（出力12桁）};
\node[nm] (PX) at (-5.4,-1.8) {Package-X 2.1.1（\texttt{LoopRefine}）\\$2+\frac{1}{\epsilon}+\mathrm{DiscB}(s;m,m)+\ln\frac{\mu^2}{m^2}$\\$\mathrm{DiscB}(2;1,1)=-\pi/2$};
\node[nm] (FH) at (5.4,-1.8) {FeynHelpers 2\\\texttt{LToolsEvaluate}：$\frac{1}{\epsilon}-1.2927418775458297$\\\texttt{PaXEvaluate}：$\frac{1}{\epsilon}+2-\gamma_E-\frac{\pi}{2}-\ln\pi$};
\draw[e] (LT) -- node[lab] {倍精度で一致} (E);
\draw[e] (CO) -- node[lab] {12桁で一致} (E);
\draw[e] (PX) -- node[lab] {有限部分が\\厳密に一致} (E);
\draw[e] (FH) -- node[lab] {差は $-(\gamma_E+\ln\pi)$\\（規格化）} (E);
\draw[e,dashed] (FH.north) -- node[lab,right] {二つの経路が\\倍精度で一致} ++(0,0.9);
\end{tikzpicture}
\caption{$B_0(2;1,1)$ の相互検算。緑が厳密値、青が Mathematica を使わない数値実装、橙が Mathematica 上の実装である。FeynHelpers の二つの経路（LoopTools 経由と Package-X 経由）は互いに一致し、厳密値とは規格化に由来する定数 $-(\gamma_E+\ln\pi)$ だけずれる。}
\label{fig:b0}
\end{figure}

\paragraph{記号計算の連携。}FeynCalc とパッチ版 FeynArts の組では、標準模型（スカラーと $W,Z$ を除く）で1ループの光子の自己エネルギーを生成し、テンソル積分を PaVe 関数まで簡約した。9種のフェルミオンのループからそれぞれ $A_0$ と $B_0$ が一つずつ現れ、計18個になった。結果は横波の構造 $p^2g^{\mu\nu}-p^\mu p^\nu$ を持ち、ゲージ不変性から期待される形になっている。パッチなし FeynArts と FormCalc の組では、$e^+e^-\to\mu^+\mu^-$ の木の振幅について \texttt{CalcFeynAmp} が FORM との連携を含めて完了した。

\paragraph{FeynRules から MadGraph へ。}FeynRules で標準模型を読み込み、UFO と FeynArts の模型を出力した。メッセージは0件だった。出力した UFO の163頂点の粒子の組を、MadGraph 同梱の \texttt{sm} の153頂点と比べると、\texttt{sm} の頂点はすべて含まれていた。余分な10頂点は、ヒッグス粒子とゴールドストーン粒子（$G^0,G^\pm$）が u・d・s クォークに結合する頂点だけである。MadGraph の \texttt{sm} は u・d・s を質量ゼロとして扱うので、これらの湯川結合を持たない。差はこれで説明できる。

\paragraph{MadGraph。}$e^+e^-\to\mu^+\mu^-$ を MadGraph の既定の設定で計算した。ビームは各 500 GeV（$\sqrt{s}=1$~TeV）、レプトンには $|\eta_\ell|<2.5$ と $p_T>10$~GeV のカットがかかる。結果は $0.10437\pm0.00018$~pb だった。同じ電弱パラメータを使い、$\gamma$ と $Z$ の交換を含む最低次の解析式で計算すると $0.10426$~pb になる（\ref{app:derivations}）。両者の差は統計誤差の0.6倍である。ループ誘起過程 $gg\to H$ は、\texttt{loop\_sm} 模型で13 TeV の LO として計算し、$15.76\pm0.05$~pb を得た（PDF は nn23lo1、$m_H=125$~GeV、$m_t=173$~GeV）。こちらは独立な検算をしておらず、MadLoop・Ninja・COLLIER の連携が動くことの確認にとどまる。

\section{使い方}\label{sec:usage}

新しいターミナルは \pth{.zshrc} を読むので、\texttt{form}、\texttt{lt}、\texttt{mg5\_aMC}、\texttt{wolframscript} がそのまま使える。MadGraph は付属の起動スクリプトを通して、conda 環境 \texttt{qft} と SDKROOT を設定した状態で起動する。

\begin{Verbatim}
form file.frm        # FORM 5.0.2 (tform for the multi-threaded version)
lt                   # LoopTools command-line front end
mg5_aMC              # MadGraph5_aMC@NLO 3.8.1 with the conda env qft and SDKROOT
\end{Verbatim}

Mathematica では、用途に応じて次のように読み込む。パッチ版とパッチなし版の FeynArts は、別々のカーネルで使う。

\begin{Verbatim}
(* FeynCalc + FeynArts + FeynHelpers *)
$LoadAddOns = {"FeynArtsLoader", "FeynHelpers"};
<< FeynCalc`

(* FeynArts + FormCalc (separate kernel) *)
<< FeynArts`
<< FormCalc`

(* Package-X and PVReduce *)
<< X`
<< X`PVReduce`

(* FeynRules; load models from the backup folder by full path *)
$FeynRulesPath = FileNameJoin[{$UserBaseDirectory, "Applications", "FeynRules"}];
SetDirectory[$FeynRulesPath]; << FeynRules`
LoadModel["~/Documents/Claude/QFTcalculation/models/backup_20261009/feynrules_Library_Models/AbelianHiggs/AbelianHiggs.fr"]
\end{Verbatim}

\paragraph{対話的な実行。}計算のスクリプトは、まとめて実行できるだけでなく、セルごとに実行して途中の式や図を確かめられるように書く。この形は $\phi^3$ 模型の手順 (5)（\nolinkurl{models/phi3/Src/scat4_feyncalc.wl} と \nolinkurl{scat4_numeric.py}）で試作した。Wolfram のスクリプト（\texttt{.wl}）には、\texttt{(* ::Section:: *)} や \texttt{(* ::Text:: *)} などのフロントエンドのセルの印を付ける。\texttt{wolframscript -file} で実行すれば、従来どおり計算のセルが順に評価され、\texttt{Print} の行だけが出力される。同じファイルを \texttt{wopen}（\ref{app:files}）で開くと、Mathematica 14.2.1 のノートブックとしてセルに分かれて表示される。ダブルクリックで開くと古い 14.1（\texttt{Wolfram.app}）が選ばれることがあるので、使わない。途中の対象を表示するセルは \texttt{If[\$Notebooks, ...]} で囲む。表示するのは、ダイアグラム、図ごとの振幅とテンソル簡約の結果、極、数値の表である。こうしたセルは、ノートブックでは「Run All Code」でも評価され、\texttt{wolframscript} では評価されない。最初は入力セル（\texttt{(* ::Input:: *)}）として書いたが、「Run All Code」はこれを評価しないので、この形に改めた。FeynArts の \texttt{Paint} の結果をそのまま評価するとラベルしか表示されないので、\texttt{Show /@ Render[Paint[..., DisplayFunction -> Identity]]} で Graphics に直す。ファイルの場所は、\texttt{\$InputFileName} が空のときは \texttt{NotebookDirectory[]} から取る。フロントエンドで保存してもこの形式は保たれ、出力のセルは保存されない。同じファイルのウィンドウが開いたままだと、\texttt{open -a} はそれを前に出すだけでファイルを読み直さない。古いウィンドウを保存すると編集が上書きされるので、保存せずに閉じてから開き直す。

Python のスクリプト（\texttt{.py}）は \texttt{\# \%\%} でセルに分け、説明は \texttt{\# \%\% [markdown]} のセルに書く。\texttt{python} で実行すれば全体が通り、図はファイルに保存される。VS Code では「Run Cell」でセルを対話ウィンドウに送ると、セルの最後の式の値と図がその場に表示される。そのため、conda 環境 \texttt{qft} に \texttt{ipykernel} 7.4.0 を加え、リポジトリの \nolinkurl{.vscode/settings.json} で \texttt{qft} の Python を指定した。対話的に実行するときは matplotlib の Agg を使わない（\texttt{"ipykernel" in sys.modules} で判定する）。\texttt{\_\_file\_\_} がないときは、作業ディレクトリ（ファイルのある場所）を使う。

\begin{Verbatim}
~/hep/bin/wopen ~/Documents/Claude/QFTcalculation/models/phi3/Src/scat4_feyncalc.wl
                               # notebook in Mathematica 14.2.1; "Run All Code" shows the Inspect cells
wolframscript -file scat4_feyncalc.wl      # batch: Inspect cells skipped, (a)/(b)/(c) lines printed
~/miniforge3/envs/qft/bin/python scat4_numeric.py   # batch; in VS Code use "Run Cell" (# %%)
\end{Verbatim}

\paragraph{全スクリプトの変換と検証。}手順 (5) で試作したあと、$\phi^3$ 模型の残りのスクリプトも同じ形にした。合わせて25本（Wolfram 15本、Python 10本）である。Inspect セルで見えるものには、たとえば次がある。ラグランジアンと Feynman 則、ダイアグラム、図ごとの振幅のテンソル簡約の前後、$D=6-2\epsilon$ での Feynman パラメータ表示、極と対抗項、$\beta$ 関数、NLOCT の対抗項、振幅の $\mu$ 微分と期待される形との差、高エネルギーでの振幅の図である。変換では計算の行を打ち直さず、セルの印と Inspect セルを決まった位置に差し込むだけにした。そのための道具と、変換が結果を変えていないことを確かめる道具を \nolinkurl{environment/interactive/} に置いた（使い方は同じ場所の \nolinkurl{README.md}）。役割は次のとおりである。
\begin{itemize}
\item \texttt{wlcells.py} と \texttt{pycells.py}：変換。
\item \texttt{nbcheck.wl}：見えないフロントエンドで Inspect セルを評価し、その出力を PDF に書き出す。
\item \texttt{pycheck.py}：IPython のカーネルで1セルずつ実行する。
\item \texttt{samefiles.sh}：書き出されたファイルをコミット済みの版と比べる。
\end{itemize}
変換したスクリプトはすべてまとめて実行し、次を確かめた。書き出された CSV はコミット済みの版と1バイト単位で一致した。図の PDF は中身の文字が一致した。確認用の出力行も変わらなかった。FeynRules と NLOCT で作り直した模型のファイルは、生成日時の行だけが変わった。Inspect セルはすべて出力を返し、Python のセルもすべて誤りなく通った。

途中で踏んだ落とし穴は四つある。第一に、フロントエンドはパッケージの形式のファイルを、セルの印と2行以上の空行でだけセルに分け、1行の空行では分けない。そのため、一つのセルに複数の文が入ることがある。たとえば \texttt{Get["FeynCalc\bq"]} と、ファイルの場所を決める文が同じセルに入る。第二に、見えないフロントエンドで評価する道具は、\texttt{wolframscript -file} ではなく、\texttt{-code} で \texttt{ReadString} から読み込む。\texttt{-file} では \texttt{\$InputFileName} が道具自身を指し、ノートブックと同じ空の状態にならないからである。セルは、\texttt{\$InputFileName} をファイルの場所に置き換えたうえで丸ごと評価する。また、道具の変数名はスクリプトの記号と衝突しないようにする。今回はスクリプトが \texttt{res} という名前を使っていて衝突した。第三に、\texttt{Paint} に \texttt{SheetHeader -> None} を付けると、\texttt{Render} が失敗する。第四に、zsh は引用符のない変数の中身を空白で分けない。そのため、複数のファイル名を一つの変数で渡すと、コピーや比較が黙って失敗する。比較の結果は、比べたファイルが実際にあることも含めて確かめる。

\paragraph{評価する前に図を見る。}2026年10月11日に、評価する Feynman 図を、評価する前に表示する取り決めを加えた。図を描く Inspect セルは、その図を評価するセルより前に置く。こうすると、Shift+Return で順に進めたとき、図を見たところで止められる。場合は三つある。
\begin{itemize}
\item FeynArts で振幅を作る Wolfram のスクリプトでは、\texttt{InsertFields} の文の直後に、図の Inspect セルを独立に置く。\texttt{CreateFeynAmp}、\texttt{TID}、\texttt{CalcFeynAmp} のセルより前である（図の数を数える \texttt{Length[CreateFeynAmp[...]]} は軽いので残してよい）。$\phi^3$ と $\phi^3+\phi^4$ の模型の13本がこの形である。
\item Python のスクリプトのうち、MadLoop と MadGraph の木の計算を使うもの（\texttt{madloop\_numeric.py} と \texttt{xsec\_numeric.py}）では、最初の \texttt{./check} の実行より前のセルで、MadGraph が \texttt{output} で描いた図（\texttt{born\_matrix.ps}、\texttt{loop\_matrix.ps}）を Ghostscript で PNG に変えて表示する。このセルは \texttt{if INTERACTIVE:} の中に置くので、まとめて実行したときの出力は変わらない。
\item 計算そのものが図を作らない場合（有限温度のノートの DRalgo、Matchotter、1ループの温度ポテンシャル）は、\nolinkurl{models/phi34_thermal/Src/diagrams_thermal.wl} が FeynArts で図を描き、\nolinkurl{figs/diagrams/} に PNG で書き出す。FeynArts は DRalgo や Matchete と同じカーネルで読み込まないので、別のスクリプトにした。評価するスクリプト（\texttt{dralgo\_check.wl}、\texttt{matchotter\_check.wl}、\texttt{nonpert\_single.py}、\texttt{thermal\_single.py}）は、最初の評価より前の Inspect セルでその PNG を表示する。
\end{itemize}
並べ替えたスクリプトをまとめて実行し、書き出したファイルがコミット済みの版と1バイト単位で一致することを確かめた。取り決めは \nolinkurl{environment/interactive/README.md} にもある。

ツールを更新したときは、試験をすべて実行し直す。MadGraph の試験を除けば数分で終わる。

\begin{Verbatim}
~/Documents/Claude/QFTcalculation/environment/tests/run_all.sh            # all ten checks
~/Documents/Claude/QFTcalculation/environment/tests/run_all.sh --no-mg5   # skip MadGraph
\end{Verbatim}

\section{熱的な道具の追加}\label{sec:thermal}

2026年10月11日に、$\phi^3+\phi^4$ 理論の有限温度のノート~\cite{phi34thermal} の検算のために、熱的ループと相転移の道具を加えた。まず、ノートの手順 (7) で、相の追跡と bounce の計算を自前の数値計算と照合するために、Python の CosmoTransitions~\cite{cosmotransitions} を入れた。次に手順 (8) で、レビューにある熱的な道具を入れた。レビューは有限温度の道具を二つの層に分けている。3次元の有効理論を作る層（DRalgo、Matchotter）と、相転移と重力波を求める層（PhaseTracer2、PT2GWFinder、BSMPT、TransitionListener）である。前者からは DRalgo~\cite{dralgo} と Matchotter~\cite{matchotter} を、後者からは PT2GWFinder~\cite{pt2gw} と、それが bounce に使う FindBounce~\cite{findbounce} を選んだ。どれも Mathematica で動き、ノートの式を直接確かめられる。同じ日のうちに、ノートの手順 (11) として、後者の残りの PhaseTracer2~\cite{phasetracer2}、BSMPT~\cite{bsmpt3}、TransitionListener~\cite{tl2} も入れた。さらに手順 (12) で、泡の壁の速さを局所熱平衡で求めるために WallGo~\cite{wallgo} を入れた。全体像を図~\ref{fig:thermal} に、版と配置を表~\ref{tab:thermal} に示す。図の色は図~\ref{fig:arch} と同じである。

\begin{figure}[hbp]
\centering
\begin{tikzpicture}
\begin{scope}[on background layer]
  \fill[panel] (0,0.35) rectangle (5.4,-4.05);
  \fill[panel] (5.65,0.35) rectangle (11.05,-4.05);
  \fill[panel] (11.3,0.35) rectangle (16.7,-4.05);
\end{scope}
\node[ptitle] at (2.7,0.35) {3次元の有効理論を作る層\\Mathematica 14.2.1};
\node[ptitle] at (8.35,0.35) {相転移と重力波を求める層\\Mathematica、C++、Python};
\node[ptitle] at (14.0,0.35) {相転移の独立な検算\\conda 環境 \texttt{qft}};
% 3D EFT
\node[tM] (DR) at (2.7,-0.95) {DRalgo 1.5.0（paclet）};
\node[tM,text width=4.1cm] (GM) at (2.95,-1.55) {GroupMath 1.1.3};
\node[tM] (MT) at (2.7,-2.5) {Matchete 0.6.0 の Matchotter\\（Applications、別のカーネル）};
\draw[cMd,thick] (DR.south west) ++(0.15,0) |- (GM.west);
% PT + GW
\node[tM] (PT) at (8.35,-0.95) {PT2GWFinder 1.1.0（paclet）};
\node[tM,text width=4.1cm] (FB) at (8.6,-1.55) {FindBounce 1.1.0（paclet）};
\node[tO] (PTR) at (8.35,-2.3) {PhaseTracer2 2.2.2（C++）};
\node[tO] (BS) at (8.35,-2.95) {BSMPT 3.2.0（C++）};
\node[tO] (TL) at (8.35,-3.6) {TransitionListener 2.1.0（環境 \texttt{tl}）};
\draw[cMd,thick] (PT.south west) ++(0.15,0) |- (FB.west);
% CosmoTransitions
\node[tO,text width=4.3cm] (CT) at (14.0,-0.95) {CosmoTransitions 2.0.7（pip）};
\node[tN,text width=4.3cm] (NP) at (14.0,-1.6) {numpy 2.5.3（\texttt{np.product} を補う）};
\node[tN,text width=4.3cm] (CTn) at (14.0,-2.35) {\texttt{tunneling1D}、\\\texttt{transitionFinder}};
\node[tO,text width=2.45cm] (WG) at (15.35,-3.45) {WallGo 1.1.2（pip）};
% target
\node[draw=cGd,fill=cGf,rounded corners=2pt,text width=16.3cm,align=flush center,font=\footnotesize,inner sep=3pt] (NOTE)
  at (8.35,-5.5) {照合の相手：$\phi^3+\phi^4$ 理論の有限温度のノート（\texttt{models/phi34\_thermal}）の、手で導いた式と自前の数値計算};
\draw[flow,cGd] (MT.south) -- node[right,font=\scriptsize,text=cGd,align=left] {温度ポテンシャルの係数\\（記号）} (MT.south |- NOTE.north);
\draw[flow,cGd] (TL.south) -- node[right,font=\scriptsize,text=cGd,align=left] {$T_c$、強さ、$T_n$、$\beta/H$、\\重力波のスペクトル} (TL.south |- NOTE.north);
\draw[flow,cGd] ([xshift=-1.45cm]CTn.south) -- node[right,font=\scriptsize,text=cGd,align=left,pos=0.62] {bounce の作用、\\$T_c$、$T_n$、相の追跡} ([xshift=-1.45cm]CTn.south |- NOTE.north);
\draw[flow,cGd] (WG.south) -- node[right,font=\scriptsize,text=cGd,align=left] {局所熱平衡の\\壁の速さ} (WG.south |- NOTE.north);
\end{tikzpicture}
\caption{熱的な道具の全体像。左の二つの層はレビューの分類による。GroupMath と Matchete は Applications に置き、DRalgo、PT2GWFinder、FindBounce は paclet として入れた。DRalgo と Matchete は同じカーネルで読み込まない。FindBounce は PT2GWFinder が bounce の計算に使う。PhaseTracer2 と BSMPT は \texttt{\textasciitilde/hep} に構築した C++ のプログラム、TransitionListener は conda 環境 \texttt{tl}、WallGo は \texttt{qft} の Python のパッケージである。緑の矢印は、ノートの何と照合したかを示す。DRalgo と Matchotter は同じ量を独立に計算した。}
\label{fig:thermal}
\end{figure}

\begin{table}[hbp]
\centering\small
\caption{熱的な道具の版、配置、入手元。Applications は表~\ref{tab:installed} と同じ、Repository は \texttt{\textasciitilde/Library/Wolfram/Paclets/Repository} を表す。}
\label{tab:thermal}
\begin{tabularx}{\linewidth}{>{\raggedright\arraybackslash}p{3.3cm}>{\raggedright\arraybackslash}p{1.1cm}>{\raggedright\arraybackslash}X>{\raggedright\arraybackslash}p{4.9cm}}
\toprule
\textbf{ツール} & \textbf{版} & \textbf{配置} & \textbf{入手元} \\
\midrule
DRalgo & 1.5.0 & Repository/DRalgo\_\_DRalgo-1.5.0 & GitHub（DR-algo）の paclet~\cite{dralgo,url:dralgo} \\
GroupMath & 1.1.3 & Applications/GroupMath & 作者のサイト~\cite{groupmath,url:groupmath} \\
Matchete（Matchotter を含む） & 0.6.0 & Applications/Matchete & GitLab の master（コミット 742c31e）~\cite{matchete,matchotter,url:matchete} \\
FindBounce & 1.1.0 & Repository/FindBounce-1.1.0 & GitHub（vguada）の paclet~\cite{findbounce,url:findbounce} \\
PT2GWFinder & 1.1.0 & Repository/PT2GW-1.1.0 & GitHub（finshky）の paclet~\cite{pt2gw,url:pt2gw} \\
CosmoTransitions & 2.0.7 & \texttt{qft} の site-packages/cosmoTransitions & PyPI~\cite{cosmotransitions,url:ct} \\
PhaseTracer2 & 2.2.2 & \texttt{\textasciitilde/hep/PhaseTracer-2.2.2} & GitHub（PhaseTracer）のタグ~\cite{phasetracer2,url:phasetracer} \\
ALGLIB（PhaseTracer2 用） & 4.08.0 & \texttt{\textasciitilde/hep/lib/libalglib.a}、\texttt{\textasciitilde/hep/include/alglib} & alglib.net の C++ のソース~\cite{url:alglib} \\
NLopt（PhaseTracer2、BSMPT 用） & 2.11.0 & Homebrew & \texttt{brew install} \\
BSMPT & 3.2.0 & \texttt{\textasciitilde/hep/BSMPT-3.2.0} & GitHub（BSMPT）のタグ~\cite{bsmpt3,url:bsmpt} \\
TransitionListener & 2.1.0 & \texttt{\textasciitilde/hep/TransitionListener-2.1.0}、conda 環境 \texttt{tl} & GitHub（tasicarl）のタグ~\cite{tl2,url:tl} \\
WallGo & 1.1.2 & \texttt{qft} の site-packages/WallGo & PyPI~\cite{wallgo,url:wallgo} \\
h5py（WallGo 用） & 3.16.0 & \texttt{qft} & conda-forge \\
\bottomrule
\end{tabularx}
\end{table}

\paragraph{導入。}どれも \texttt{curl}、\texttt{git}、\texttt{pip} で取得し、ブラウザは使わなかった（第~\ref{sec:issues} 節）。DRalgo、FindBounce、PT2GWFinder は、GitHub で配布されている paclet を \texttt{PacletInstall} で入れた。中身は Mathematica のパッケージと文書で、コンパイル済みのライブラリは含まない。GroupMath は作者のサイトの zip を Applications に展開した。DRalgo は読み込むときに GroupMath を Applications から読む。Matchete には版の付いた配布がない。そこで GitLab の master（2026年9月24日のコミット 742c31e、版の表示は 0.6.0）を取得し、\texttt{.git} を除いて Applications に写した。これは Matchete の \texttt{install.m} が行う操作と同じである。版の番号だけでは中身が決まらないので、入れ直すときはこのコミットを指定する。CosmoTransitions は PyPI の wheel を \texttt{pip install --no-deps} で \texttt{qft} に入れた。必要な numpy と scipy はすでに入っている。大きさは Matchete が 90 MB、PT2GWFinder が 16 MB で、ほかは 5 MB 以下である。コマンドは \nolinkurl{files/build_commands.txt} の8番目の項目にある。PhaseTracer2 が使う ALGLIB は Homebrew にないので、alglib.net の C++ のソースを静的ライブラリにして \pth{hep/lib} に置いた。NLopt は Homebrew で入れた。BSMPT は conan を使わず、Homebrew の Eigen、GSL、nlohmann-json、NLopt で構築した。Homebrew の GSL には CMake の設定ファイルがないので、CMake の \texttt{FindGSL} を呼ぶ1行のファイルを \nolinkurl{BSMPT-3.2.0/cmake-shim} に置いた。libcmaes、重粒子の生成の計算、テストは外した。TransitionListener は、多くの依存（ultranest、mpi4py、getdist など）を \texttt{qft} に入れないよう、別の conda 環境 \texttt{tl} を作って入れた。構築のスクリプトとログは \pth{hep/src} にあり、コマンドは \nolinkurl{files/build_commands.txt} の9番目の項目にある。WallGo は h5py、numpy、scipy を要る。\texttt{qft} になかった h5py（3.16.0）を conda-forge から入れ（試しの実行で、既存のパッケージは変わらず新しいものが加わるだけと確かめた）、PyPI の wheel（\pth{hep/src/wallgo-dl} に保存）を \texttt{pip install --no-deps} で入れた。平衡からずれた粒子の摩擦を求める補助のモジュール（WallGoMatrix、WallGoCollision）は入れていない。コマンドは \nolinkurl{files/build_commands.txt} の10番目の項目にある。

\paragraph{検証。}照合はノートの手順 (7)、(8) として行った。結果を表~\ref{tab:thermaltests} にまとめる。基準は、ノートで手で導いた式と、ノートの自前の数値計算（bounce の撃ち込み法と相の追跡）である。DRalgo の2ループの熱質量は、単一場の模型（$\alpha=0.9$、$\kappa_4=1$）のスピノーダル温度で1ループの熱質量 $\lambda T^2/24$ の $0.86\%$ で、最低次のスピノーダル温度を $0.43\%$ 下げた。PhaseTracer2、BSMPT、TransitionListener との照合は、ノートの手順 (11) として、一般形の5点で行った（表の8〜12）。WallGo による局所熱平衡の壁の速さは、ノートの手順 (12) として求めた（表の13、14）。細部はノートの第7節と付録~H にある。

\begin{table}[hbp]
\centering\small
\caption{熱的な道具の検証の結果。「期待値」はノートで独立に求めた値である。一般形は Dine らの形の温度ポテンシャルで、SM 的な係数と BP1〜BP4 の5点を使った。比は道具の値をノートの値で割ったもの、差は相対差である。}
\label{tab:thermaltests}
\begin{tabularx}{\linewidth}{>{\raggedright\arraybackslash}p{0.35cm}>{\raggedright\arraybackslash}p{2.5cm}>{\raggedright\arraybackslash}X>{\raggedright\arraybackslash}X>{\raggedright\arraybackslash}p{1.9cm}}
\toprule
& \textbf{対象} & \textbf{期待値} & \textbf{結果} & \textbf{判定} \\
\midrule
1 & DRalgo 1.5.0 & 4次元の $\beta$ 関数が零温度のノートと同じ。3次元の有効理論から組み直した1ループのポテンシャルが、ノートの温度ポテンシャルと同じ & $\beta$ 関数は同じ。場の微分の差は恒等的に0 & 記号で一致 \\
2 & Matchotter（Matchete 0.6.0） & DRalgo と同じ3次元の係数 & 差は0（対数 $\ln(\bar\Lambda^2/\Lambda_T^2)$ を $L_b$ と読む） & 記号で一致 \\
3 & FindBounce 1.1.0 & 木のポテンシャルの $S_3$（$\alpha=0.9$） & 1001点で差 $2\times10^{-6}$ & 一致 \\
4 & PT2GWFinder 1.1.0 & 一般形の $T_c$、強さ $\alpha_\theta$、$T_n$、$\beta/H$ & 差は $8\times10^{-7}$、$10^{-11}$、$2.4\times10^{-4}$、$5\times10^{-4}$ 以内（BP4 の $\beta/H$ だけ $+1.6\%$） & 一致（BP4 の $\beta/H$ は作用のばらつきの範囲） \\
5 & PT2GWFinder の音波 & テンプレートの違いから、全体の量の比 $1.42$ & 4点とも $1.419$ & 一致 \\
6 & CosmoTransitions 2.0.7 & bounce の作用（19通り）、一般形の $T_c$、$T_n$ & 差は $9\times10^{-8}$、$2\times10^{-15}$、$6\times10^{-9}$ 以内 & 一致 \\
7 & CosmoTransitions（単一場） & 大域的な最小の相が途切れず、臨界温度も転移もない & 臨界温度0個、転移なし & 期待どおり \\
8 & PhaseTracer2 2.2.2 & 一般形の $T_c$、$S_3/T=140$ の温度、強さ、作用（許容誤差 $10^{-10}$） & 差は $2\times10^{-15}$、$8\times10^{-7}$、$1.2\times10^{-8}$、$9\times10^{-8}$ 以内 & 一致 \\
9 & PhaseTracer2 の $\beta/H$ と重力波 & BP2〜BP4 の値 & 差は $10^{-3}$、$2\times10^{-3}$ 以内（BP1 は $+1.9\%$、SM 的な係数は $+20\%$） & 一致（速い転移では差分の幅が広すぎる） \\
10 & BSMPT 3.2.0 & 一般形の $T_c$、$T_n$（BSMPT の $H$）、$\beta/H$ & 差は $3\times10^{-8}$、$3\times10^{-5}$、$3.5\times10^{-4}$ 以内 & 一致 \\
11 & BSMPT の強さ & ノートの $\alpha_\theta$（$g_*=106.75$） & 5点とも $+1.5$〜$1.6\%$ & 説明できる（SM の $g_*(T)$ の表） \\
12 & TransitionListener 2.1.0 & BP1〜BP4 の $T_c$、浸透温度での強さと $\beta/H$ & 差は $3\times10^{-13}$、$8\times10^{-7}$、$3\times10^{-4}$ 以内 & 一致 \\
13 & WallGo 1.1.2 & 浸透温度での局所熱平衡の壁の速さ（音速一定の近似）が、TransitionListener の整合の式に相に沿った音速を入れた値と同じ & BP1〜BP3 で差は $10^{-6}$ 以内、BP4 はどちらも解なし。状態方程式の温度依存を含めると $3\times10^{-5}$〜$3\times10^{-3}$ 大きい & 一致 \\
14 & TransitionListener の壁の速さ & 整合の式に場を固めた音速を入れた値 & 差は $2\times10^{-7}$ 以内。WallGo より $0.1$〜$0.5\%$ 遅く、BP4 では存在しない解 $0.797$ を返していた & 説明できる（音速の取り決め） \\
\bottomrule
\end{tabularx}
\end{table}

\paragraph{注意。}照合の途中で、次のことが分かった。
\begin{itemize}
\item DRalgo と Matchete は同じカーネルで読み込まない。スクリプトも分けた。Matchete は、フロントエンドのない \texttt{wolframscript} ではロゴを描けずに ``Loading failed!'' と出すが、計算には影響しない。
\item Matchotter は $\hbar=1/(4\pi)^2$、$\Lambda_T=4\pi T$ と書き、$1/\epsilon$ の極を残す。その対数 $\ln(\bar\Lambda^2/\Lambda_T^2)$ が DRalgo の $L_b=2\ln(\bar\mu/T)+2\gamma_E-2\ln4\pi$ に等しいのは、$\bar\Lambda=\bar\mu e^{\gamma_E}$ のときである。この取り決めは、論文の付録 B の式 (B.4)、(B.5)（和積分の因子 $(\bar\Lambda^2C)^\epsilon$、$C=e^{-\gamma_E}/(4\pi)$）と一致する。論文は $\bar\Lambda$ を $\overline{\mathrm{MS}}$ の尺度と呼ぶので、$\bar\Lambda=\bar\mu$ とすると対数が $2\gamma_E$ ずれる。
\item DRalgo の NLO の質量と線形項は、1ループと2ループの項を合わせて返す。2ループの項は、$T^2$ に比例する項と、3次元の結合（\texttt{lam3d}、\texttt{gg3d}、\texttt{Log[\textbackslash[Mu]3/\textbackslash[Mu]]}）を含む項である。
\item FindBounce は、偽の真空から与えた端の点までの直線の経路で bounce を求める。端を bounce の中心より遠い深い真の極小に置くと、場の点を増やすほど正しい値から離れた（$\alpha=0.5$ の木のポテンシャルで、31点で $-0.42\%$、1001点で $-2.1\%$）。端を $f(0)+0.3(f_--f(0))$ に置くと、1001点で $-5\times10^{-4}$ だった。PT2GWFinder の \texttt{Action} は、端を抜け出しの点に移して計算し直すので、この問題を避けている。一般形の $S_3/T$ の差は、点の数51で $5\times10^{-4}$〜$6\times10^{-3}$、801で $5\times10^{-6}$〜$8\times10^{-5}$ だった。
\item PT2GWFinder の \texttt{FindNucleation} は、二分法の許容誤差が探す範囲の $10^{-3}$ なので、範囲を狭めて使う。\texttt{SearchPhases} は、核生成を $\int(dT/T)\,\Gamma/H^4=1$ で決め、重力波を浸透温度で評価する。ハッブル率には $|\Delta V|$ を含め、換算 Planck 質量は $2.4353\times10^{18}$~GeV である。相を純関数で渡すと、\texttt{FindCritical} が定義域を \texttt{FunctionDomain} で決めて失敗するので、\texttt{TracePhases} で作った相を渡す。音波の重力波は LISA の作業部会の2024年の形~\cite{caprini24} で、全体の量はノートが使う Hindmarsh らの規格化の1.42倍になる。batch で動かすと進み具合の表示が行を上書きするので、結果はファイルから読む。
\item CosmoTransitions は numpy 2 で消えた \texttt{np.product} を使うので、読み込む前に \texttt{np.product = np.prod} とする。\texttt{SingleFieldInstanton.findProfile} の既定の許容誤差では、作用の相対誤差が $6\times10^{-4}$ ほどある。\texttt{xtol}$=$\texttt{phitol}$=10^{-10}$、\texttt{npoints}$=4000$ とすると $10^{-7}$ 以下になる。\texttt{fullTunneling} は、経路に沿ったポテンシャルを既定で100点のスプラインで近似する（\texttt{V\_spline\_samples=None} で直接評価する）。相の追跡の刻みは既定で温度の範囲の $10^{-3}$ 倍で、狭い温度の幅を飛び越えることがある（\texttt{dtstart} と \texttt{tjump} を小さくする）。自作の核生成の基準 \texttt{nuclCriterion} は、作用が0のとき負を返さなければならない。\texttt{findAllTransitions} は準安定な相が終わる温度で止まるので、その少し下から \texttt{tunnelFromPhase} を呼ぶ。
\item PhaseTracer2 は核生成を $S_3/T=140$ で決める。$\beta/H$ は前後 $\pm0.1$~GeV の6点の作用に当てはめた直線から求めるので、非常に速い転移（$\beta/H\gtrsim10^4$）ではずれる。撃ち込み法の既定の許容誤差 $10^{-4}$ では作用に $10^{-3}$ 程度の誤差が出る（\texttt{PD\_xtol}、\texttt{PD\_phitol} を小さくする）。相の探索では、$V''(0)\to0$ となる温度より上に下限を置き、ヘッセ行列の特異性の判定を切り、場の範囲の既定の上限 $1600$~GeV を広げる。ALGLIB 4.08 では、付属の単体テスト13件のうち2件が $0.16\%$、$0.76\%$ のずれで通らない。
\item BSMPT の模型は C++ のクラスで、任意のポテンシャルは \texttt{VEff} を書き換えて与える（\texttt{diff}$=-1$ は温度微分で、強さに使う）。破れた相は \texttt{vevTreeMin} から追うので、探索の下限の温度での最小をそこに入れる。重力波の段階は SM の粒子の質量（壁の圧力）を問い合わせるので、粒子を持たないポテンシャルでは質量の表を0で初期化し、重力波の値は使わない。$g_*(T)$ は SM の表（$160$〜$215$~GeV で $105.0$〜$105.2$）、換算 Planck 質量は $2.4\times10^{18}$~GeV である。付属の \nolinkurl{standalone/} の例はテスト用のライブラリを要るので、テストを切った構築では作られない。
\item TransitionListener では、任意のポテンシャルを \texttt{generic\_potential} の派生クラスで与え、\texttt{Vtot}、エネルギー密度、放射のエネルギー密度と \texttt{V1T\_from\_X}（温度微分した勾配がこれを微分する）を書き換える。報告される $T_n$ は、泡の数の積分 $N(T)=1$（$v_w$ は局所熱平衡の模型から求める）で決まり、$\Gamma=H^4$ ではない。非常に速い転移（$\beta/H\sim10^6$）では、浸透の計算が止まる。局所熱平衡の壁の速さ（\texttt{calcWallVelocityLTE}）は、音速の2乗の温度の2階微分を、場を固めて差分で求める。相に沿った2階微分は $\partial_T^2V-(\partial_\phi\partial_TV)^2/\partial_\phi^2V$ なので、第2項が落ちて音速が大きく出る（一般形の BP1 で $0.332$ と $0.309$、BP4 で $0.406$ と $0.362$）。
\item WallGo では、任意のポテンシャルを \texttt{EffectivePotential} の派生クラスで与え、$\phi$ によらない放射の項 $-\pi^2g_*T^4/90$ を含める（状態方程式を作るため）。既定の数値の極小（Nelder--Mead）と差分では、音速一定の近似の $c_b^2$ が $10\%$ ほどずれることがあり（一般形の BP1 の浸透温度で $0.340$、厳密には $0.309$）、\texttt{epsilon} が負だとして計算が止まる。多項式なら、\texttt{findLocalMinimum}、\texttt{derivT}、\texttt{derivField}、\texttt{deriv2FieldT}、\texttt{deriv2Field2}、\texttt{allSecondDerivatives} を厳密な式で書き換える。\texttt{derivT} は1点ではスカラーを返すようにする（配列を返すと、相の追跡で配列の次元が合わずに止まる）。相を追う刻みは \texttt{temperatureVariationScale} の $10^{-6/4}$ 倍である。大きすぎると補間が粗くなり、小さすぎると、破れた相が終わる温度の近くで求めた Jouguet の速さが下がり、根を見落として $1$ を返す。尺度を2倍と半分にして値を確かめる。\texttt{wallSpeedLTE} はデトネーションを返さず、局所熱平衡のデフラグレーションも hybrid もなければ $1$ を返す。
\end{itemize}

\paragraph{使い方。}読み込みは次のとおりである。照合は、ノートの \nolinkurl{Src/} にある四つの実行スクリプトで実行し直す。Mathematica の3本は合わせて約40分、PhaseTracer2、BSMPT、TransitionListener は約5分、WallGo は約2分かかる。\texttt{tests/run\_all.sh} に \texttt{--thermal} を付けると、これらもまとめて実行する。

\begin{Verbatim}
(* DRalgo; GroupMath is read from Applications *)
DRalgo`DRalgo`$LoadGroupMath = True;
<< DRalgo`DRalgo`
(* Matchete / Matchotter: a kernel of its own *)
<< Matchete`
(* PT2GWFinder; loads FindBounce *)
<< PT2GW`
$PT2GWPrint = False;
\end{Verbatim}
\begin{Verbatim}
# CosmoTransitions (conda env qft)
import numpy as np; np.product = np.prod      # numpy 2 removed np.product
from cosmoTransitions import tunneling1D, generic_potential, transitionFinder
# WallGo (conda env qft); the model as in models/phi34_thermal/Src/wallgo_check.py
import WallGo
\end{Verbatim}
\begin{Verbatim}
cd ~/Documents/Claude/QFTcalculation/models/phi34_thermal/Src
./run_cosmotransitions_check.sh
./run_thermal_tools_check.sh        # DRalgo, Matchotter, PT2GWFinder (about 40 minutes)
./run_review_tools_check.sh         # PhaseTracer2, BSMPT (compiled here), TransitionListener (env tl)
./run_wallgo_check.sh               # WallGo (env qft), TransitionListener's LTE matching (env tl)
~/Documents/Claude/QFTcalculation/environment/tests/run_all.sh --no-mg5 --thermal   # everything
\end{Verbatim}

\section{2ループの道具の追加}\label{sec:twoloop}

2026年10月12日に、$\phi^3$ 理論のノート~\cite{phi3note} の手順 (16)（$d=4$ の2ループの自己エネルギー）のために、レビューが推奨した2ループの道具を加えた。IBP 簡約の Kira 3.1~\cite{kira3}（代数のソルバ Fermat 7.9b~\cite{url:fermat} とともに）と FIRE 7.1~\cite{fire7}、マスター積分を数値で求める AMFlow 2.0~\cite{amflow,amflow2}、pySecDec 1.6.6~\cite{pysecdec}、feyntrop 1.1.0~\cite{feyntrop}、2ループの自己エネルギーの基底積分のライブラリ TSIL 1.46~\cite{tsil} である。FeynCalc の積分族は、FeynHelpers 2 を通して Kira と FIRE に渡した。版と配置を表~\ref{tab:twoloop} に示す。

\begin{table}[hbp]
\centering\small
\caption{2ループの道具の版、配置、入手元。}
\label{tab:twoloop}
\begin{tabularx}{\linewidth}{>{\raggedright\arraybackslash}p{3.4cm}>{\raggedright\arraybackslash}p{1.2cm}>{\raggedright\arraybackslash}X>{\raggedright\arraybackslash}p{4.2cm}}
\toprule
\textbf{ツール} & \textbf{版} & \textbf{配置} & \textbf{入手元} \\
\midrule
Kira & 3.1 & \texttt{\textasciitilde/hep/Kira-3.1}（FireFly、FLINT 2.8.4、yaml-cpp を同梱で構築）、起動スクリプト \texttt{\textasciitilde/hep/bin/kira} & GitLab（kira-pyred）のタグ~\cite{kira3,url:kira} \\
Fermat（Kira 用） & 7.9b & \texttt{\textasciitilde/hep/fermat-7.9b}（Apple M 版） & 作者のサイト~\cite{url:fermat} \\
CLN、GiNaC（Kira 用） & 1.3.7、1.8.10 & \texttt{\textasciitilde/hep/lib}（静的）、\texttt{\textasciitilde/hep/bin/ginsh} & ginac.de~\cite{url:ginac} \\
FIRE & 7.1 & \texttt{\textasciitilde/hep/FIRE-7.1} $\to$ \texttt{\textasciitilde/hep/src/fire-7.1}（\texttt{FIRE7/bin}、\texttt{FIRE7.m}） & MSU の GitLab のタグ~\cite{fire7,url:fire} \\
AMFlow & 2.0 & \texttt{\textasciitilde/hep/AMFlow-2.0}（C++ の微分方程式のソルバ \texttt{desolver}、\texttt{link}） & GitLab（multiloop-pku）のタグ~\cite{amflow2,url:amflow} \\
MPSolve、yaml-cpp（AMFlow 用） & 3.2.3、0.8.0 & \texttt{\textasciitilde/hep/bin/mpsolve}、\texttt{\textasciitilde/hep/lib/libyaml-cpp.a} & 作者のサイト、GitHub \\
pySecDec & 1.6.6 & \texttt{qft} の site-packages（ソースから構築） & PyPI~\cite{pysecdec,url:pysecdec} \\
sympy、nest-asyncio（pySecDec 用） & 1.14.0、1.6.0 & \texttt{qft}（gmpy2 2.3.2、gmp、mpfr、mpc が加わった） & conda-forge \\
GNU make（pySecDec 用） & 4.4.1 & Homebrew（\texttt{gmake}） & \texttt{brew install make} \\
feyntrop & 1.1.0 & \texttt{\textasciitilde/hep/feyntrop-1.1.0} & GitHub（michibo）のタグ~\cite{feyntrop,url:feyntrop} \\
TSIL & 1.46 & \texttt{\textasciitilde/hep/lib/libtsil.a}、\texttt{\textasciitilde/hep/include/tsil.h}（試験は \texttt{\textasciitilde/hep/src/tsil-1.46}） & 作者のサイト~\cite{tsil,url:tsil} \\
meson、ninja（構築用） & 1.12.1、1.13.2 & 仮想環境 \texttt{\textasciitilde/hep/opt/buildtools} & PyPI \\
\bottomrule
\end{tabularx}
\end{table}

\paragraph{導入。}どれも \texttt{curl}、\texttt{git}、\texttt{pip} で取得した。Homebrew で新しく入れたのは GNU make だけである。Homebrew の GiNaC と meson は、python@3.15 と openssl@4 を入れ、python@3.14、readline、sqlite、xz、ca-certificates を更新しようとした。そこで CLN と GiNaC はソースから静的に \pth{hep} に構築し、meson と ninja は専用の仮想環境に入れた。Kira は meson で構築し、FireFly、FLINT、yaml-cpp を同梱の副計画として作った。FLINT の副計画は Homebrew の gmp の場所を知らないので、\texttt{LIBRARY\_PATH} を与えた。Fermat はフリーウェアで、機関や研究費のある利用者には1台あたり US\$75 の寄付が求められ、営利目的では使えない。寄付は利用者が判断する。FIRE は、README のとおり Homebrew の LLVM で構築した。Homebrew の clang は macOS の SDK を見つけないので、\texttt{-isysroot} を与える小さなスクリプトを置いた（\ref{app:files}）。その C++ の標準ライブラリのヘッダはシステムの libc++ にない記号を使うので、リンクでは Homebrew の libc++ をつないだ。FIRE の Makefile は \texttt{-Werror} で gcc 向けの警告を使うので、FIRE 用のスクリプトで \texttt{-Wno-error} を後ろに付けた。FIRE の式の簡約器 fuel が要る FLINT は、記録されたコミット（49b15d9、3.1.0-dev）を CMake で BLAS なしに構築した。Homebrew の FLINT 3.6 とは API が合わない。fuel の部分モジュールのうち、使わない pari と symbolica は取得しなかった（PARI の取得は非常に遅い）。zstd と snappy の圧縮は構築していない。FIRE は kyotocabinet を絶対パスでつなぐので、ソースのフォルダを動かさず、\pth{hep/FIRE-7.1} からリンクを張った。AMFlow は Mathematica のパッケージで、\texttt{ibp\_interface} の Kira と FIRE のパスを書き換えた（元のファイルは \texttt{.orig}）。C++ の微分方程式のソルバは、付属の \texttt{auto\_install.sh}（sudo を使い、\pth{.zshrc} を書き換える）を使わず、MPSolve と yaml-cpp を \pth{hep} に構築してから \texttt{make} で作った。pySecDec は PyPI のソースを \texttt{pip install --no-deps} で \texttt{qft} に入れた。依存の sympy を pip で入れると mpmath 1.4.1 が 1.3.0 に下げられるので、conda-forge から入れた。このとき gmpy2 が加わり、mpmath の下回りが gmpy に替わったので、mpmath を使う $\phi^3$ と $\phi^3+\phi^4$ のスクリプト8本を実行し直し、CSV が1バイト単位で、図が中身の文字で変わらないことを確かめた。pySecDec が生成する Makefile は GNU make 4 以上を要る（macOS の make は 3.81）。feyntrop は Homebrew の LLVM の OpenMP で構築した。TSIL は倍精度で構築した。Apple silicon では long double は double と同じなので、試験の基準も倍精度のものを使う。コマンドは \nolinkurl{files/build_commands.txt} の11番目の項目にある。

\paragraph{検証。}各ツールの付属の試験と例を通し、$\phi^3$ のノートの手順 (16) で互いに突き合わせた。結果を表~\ref{tab:twolooptests} にまとめる。手順 (16) では、Kira と FIRE の簡約が記号として一致し、8個のマスター積分が四つの方法（AMFlow、TSIL、pySecDec、feyntrop）で一致し、自己エネルギーの有限部分が文献~\cite{tsil}の式と相対 $8\times10^{-14}$ 以内で一致した。細部はノートの第8節と付録~T、U にある。

\begin{table}[hbp]
\centering\small
\caption{2ループの道具の検証の結果。}
\label{tab:twolooptests}
\begin{tabularx}{\linewidth}{>{\raggedright\arraybackslash}p{0.35cm}>{\raggedright\arraybackslash}p{2.5cm}>{\raggedright\arraybackslash}X>{\raggedright\arraybackslash}X>{\raggedright\arraybackslash}p{1.9cm}}
\toprule
& \textbf{対象} & \textbf{期待値} & \textbf{結果} & \textbf{判定} \\
\midrule
1 & Kira 3.1 & 付属の試験（FireFly 5件、Kira 2件を参照の結果と照合） & \texttt{-O2} の構築で7件とも通る。\texttt{-O3} では Kira の2件が落ち、すべての積分が0に簡約される & 一致（\texttt{-O2}） \\
2 & Kira 3.1 & 1ループの質量のある泡と tadpole の簡約 & 係数が解析式と一致 & 一致 \\
3 & FIRE 7.1 & 付属の試験（\texttt{tests.makefile}、FLINT） & 17件の照合がすべて一致（zlib、zstd、snappy は飛ばす） & 一致 \\
4 & AMFlow 2.0 & ソルバの確認、付属の例を Kira と FIRE で解いた値と付属の結果 & すべての係数が10桁で一致 & 一致 \\
5 & pySecDec 1.6.6 & $\int(x+y)^{-2+\epsilon}=1/\epsilon+1-\ln2+O(\epsilon)$ & 差 $5\times10^{-15}$ & 一致 \\
6 & feyntrop 1.1.0 & 付属の試験（7例の参照値） & 比はすべて1 & 一致 \\
7 & TSIL 1.46 & 付属の試験（320件） & 通過296、警告23、失敗1（作者の記録では倍精度で303、16、1） & 期待どおり \\
8 & Kira と FIRE & $\phi^3$ の2ループの自己エネルギーの3つの積分の簡約 & 対称性で等しい積分をそろえると記号として一致 & 一致 \\
9 & AMFlow、TSIL、pySecDec、feyntrop & 8個のマスター積分（4点） & TSIL $7.5\times10^{-14}$、pySecDec $2\times10^{-9}$、feyntrop は誤差の範囲 & 一致 \\
10 & 自己エネルギー & 極が消える。有限部分は文献~\cite{tsil}の $-(M+V)/2$ & 極は AMFlow の30桁の値で29桁以上まで0、有限部分は $8\times10^{-14}$ 以内 & 一致 \\
\bottomrule
\end{tabularx}
\end{table}

\paragraph{注意。}導入と照合の途中で、次のことが分かった。
\begin{itemize}
\item Kira 3.1 を meson の既定（\texttt{buildtype=release}、\texttt{-O3}）で Apple clang 16 で構築すると、すべての積分を0に簡約する。メッセージは出ず、付属の試験で分かる。デバッグの構築（\texttt{-O0}）と \texttt{-O2}（\texttt{debugoptimized}）は正しい。FLINT を外しても変わらない。
\item Kira はマスター積分の候補を種の範囲で選ぶ。FeynHelpers の \texttt{KiraCreateJobFile} が目標の積分から決める種（$r=5$、$s=0$）では、日没型のマスター積分が一つ多く残った。分子を含む種（$s=2$）で FIRE と同じ数に減った。優先するマスター積分は job ファイルの \texttt{preferred\_masters} にファイルで与える。同じフォルダに前回のデータベースがあると、Kira は「すでに簡約済み」として計算をやり直さない。
\item FIRE は対称性で等しい積分を別々のマスター積分として残すので、比べる前にそろえる（FeynCalc の \texttt{FCLoopFindIntegralMappings}）。伝播関数2本の積分は正しく扱えないと FIRE 自身が警告する（issue 40）。FeynHelpers の既定の圧縮 \texttt{zstd} は、zstd なしで構築した FIRE では使えないので \texttt{lz4} にする。優先するマスター積分は \texttt{\#preferred} で与える。
\item AMFlow の \texttt{SolveIntegrals[ints, goal, order]} の \texttt{order} は、$L$ ループの族の $\epsilon^{-2L}$ から数えた次数である。2ループの族に2を与えると $1/\epsilon^2$ の項しか求めず、残りの係数は警告なしに0になる。同じ族の名前で運動学を変えて呼んでも、結果は正しかった。付属の \texttt{check\_install.sh} は古い Mathematica 14.1（\texttt{Wolfram.app}）を使い、計算が終わってもカーネルが終了しない。既定の C++ のソルバ（WSTP の \texttt{link}）では、$\phi^3$ の照合の3回の実行のうち2回が、小さな1ループの部分系で止まった。並列のサブカーネルの一つが CPU を使い続け、\texttt{link} は待ったままで、ほかの計算を走らせていなくても起きた。照合では Mathematica のソルバ（\texttt{SetAMFOptions["DESolver" -> "MMA"]}）を使う。AMFlow はキャッシュを実行中のファイルの隣に作る（バッチでは \texttt{\$InputFileName}、ノートブックでは \texttt{NotebookFileName[]} の場所）。セルをノートブックを通さずに評価すると後者が \texttt{\$Failed} になって止まるので、照合のスクリプトでは \texttt{CacheDirectory} を定義し直して作業フォルダに向けた。
\item pySecDec のパッケージの名前に \texttt{I} は使えない（sympy の虚数単位）。生成したパッケージは \texttt{gmake} で構築する。積分は $\int\prod\mathrm{d}^Dk/(i\pi^{D/2})\prod(q^2-m^2+i0)^{-\nu}$ の規格化で返る。
\item feyntrop の Python の部分は、作業フォルダの \texttt{./feyntrop} を呼ぶ。規格化はユークリッドの測度 $\mathrm{d}^Dk/\pi^{D/2}$ に $\Gamma$ の前因子を付けたもので、TSIL の太字の積分と同じである。
\item TSIL の \texttt{tsil.h} は \texttt{complex.h} を読むので、利用者のプログラムで \texttt{I} を変数名に使えない。太字の真空の日没型 $I$ は直接は返らない。
\item FIRE と feyntrop は Homebrew の LLVM の libc++ と libomp につないでいる。Homebrew で LLVM を更新したら、構築し直して試験を通す。
\end{itemize}

\paragraph{使い方。}FeynCalc からの使い方は次のとおりである。照合はノートの \nolinkurl{Src/run_twoloop4_check.sh} で実行し直す（約12分）。\texttt{tests/run\_all.sh} に \texttt{--twoloop} を付けると、これも実行する。
\begin{Verbatim}
(* FeynHelpers 2 -> Kira 3.1 and FIRE 7.1 (see models/phi3/Src/twoloop4_feyncalc.wl) *)
KiraCreateIntegralFile[ints, topo, dir];  KiraCreateConfigFiles[topo, ints, dir, KiraMassDimensions -> {m -> 1, pp -> 2}];
(* write dir/<topo>/job.yaml with r, s and preferred_masters, then *)
KiraRunReduction[FileNameJoin[{dir, "<topo>"}], KiraBinaryPath -> "~/hep/Kira-3.1/bin/kira",
  KiraFermatPath -> "~/hep/fermat-7.9b/fer64"];  KiraImportResults["<topo>", dir]
FIREPrepareStartFile[topo, dir, FIREPath -> "~/hep/FIRE-7.1/FIRE7/FIRE7.m"];
FIRECreateStartFile[dir, topo, FIREMathematicaKernelPath -> "~/hep/opt/Wolfram142/Contents/MacOS/WolframKernel"];
FIRECreateConfigFile[topo, dir, FIRECompressor -> "lz4"];
FIRERunReduction[FileNameJoin[{dir, "<topo>"}], FIREBinaryPath -> "~/hep/FIRE-7.1/FIRE7/bin/FIRE7"]
(* AMFlow 2.0, own kernel *)
Get["~/hep/AMFlow-2.0/AMFlow.m"]; SetReductionOptions["IBPReducer" -> "Kira"];
SolveIntegrals[{j[fam, ...]}, 30, 2 L + 2]          (* digits, orders counted from eps^(-2L) *)
\end{Verbatim}
\begin{Verbatim}
cc -DTSIL_SIZE_DOUBLE -I$HOME/hep/include prog.c -L$HOME/hep/lib -ltsil -lm      # TSIL
~/miniforge3/envs/qft/bin/python generate.py && gmake -C name disteval -j8       # pySecDec
cd ~/Documents/Claude/QFTcalculation/models/phi3/Src && ./run_twoloop4_check.sh
\end{Verbatim}

\section{議論と今後}\label{sec:discussion}

\paragraph{検証していないこと。}今回の検証は、各ツールが起動し、代表的な計算で正しい値を返すことの確認である。次の機能は確かめていない。FeynRules の NLOCT による1ループの対抗項と NLO 用 UFO、FormCalc が生成する数値計算用 Fortran コードのコンパイルと実行、MadGraph の NLO QCD 計算（\texttt{[QCD]}）、Mathematica 14.1.0 での動作である。また、MadGraph の $gg\to H$ には独立な検算がない。表~\ref{tab:tests} の LoopTools の $C_0$ の検算は外運動量がゼロの点だけで、外運動量がゼロでない場合の誤りはその後に見つかった（第~\ref{sec:issues} 節）。これらを使う計算の前には、既知の結果の再現を最初に確かめるべきである。熱的な道具（第~\ref{sec:thermal} 節）は、1成分の実スカラーの模型でだけ確かめた。ゲージ場やフェルミオンを含む模型、DRalgo の2ループの有効ポテンシャル、多成分の場の相転移、Matchete による零温度の EFT の matching は試していない。WallGo は局所熱平衡の壁の速さだけを確かめ、平衡からずれた粒子を含む壁の速さの計算は試していない。2ループの道具（第~\ref{sec:twoloop} 節）は、質量が一つの2ループの自己エネルギーでだけ確かめた。多くの質量や外線を持つ積分、FIRE の MPI と多素数のモード、Kira の FireFly による再構成、pySecDec の領域展開、AMFlow の高ループの積分は試していない。

\paragraph{保守上の注意。}FeynRules 2.3.49 は2021年から更新されておらず、\texttt{FeynRules.m} と \texttt{Core/Loop.m} の二つの局所パッチを前提に動いている。FeynRules を入れ直したときは、\ref{app:files} の差分を当て直す必要がある。Package-X は開発が止まっており、$O(\epsilon)$ の項を持たない。FeynCalc の公式インストーラは、タグではなく安定版の枝（\texttt{hotfix-stable}）を取得する。今回入ったのは 10.2.1 だが、次に実行したときに同じ版になるとは限らない。FORM は 5.0.0 と 5.0.1 に深刻なバグがあったので、5.0.2 より古い版には戻さない。macOS や Homebrew の gcc を更新すると、SDKROOT や gfortran の挙動が変わりうる（gcc は 14.2 と 15.2 の二つが入っている）。いずれの場合も、更新のあとに \texttt{run\_all.sh} を実行して表~\ref{tab:tests} の値を再現できることを確かめる。熱的な道具では、Matchete に版の付いた配布がないので、入れ直すときは記録したコミットを指定する。CosmoTransitions は \texttt{np.product} を補って numpy 2 で使っているので、numpy を更新したときは第~\ref{sec:thermal} 節の照合のスクリプトを実行し直す。PhaseTracer2 と BSMPT は Homebrew の Boost、Eigen、GSL、NLopt の共有ライブラリにつないでいるので、これらを更新したら \nolinkurl{run_review_tools_check.sh} を実行し直す。Kira は \texttt{-O2} で構築する（\texttt{-O3} は誤る）。FIRE と feyntrop は Homebrew の LLVM の libc++ と libomp に、FIRE の fuel は記録したコミットの FLINT につないでいるので、LLVM を更新したら構築し直して \nolinkurl{run_twoloop4_check.sh} を実行する。

\paragraph{今後。}レビューが推奨した構成のうち、EFT・暗黒物質用のツール（DsixTools、wilson、DirectDM、MadDM）は、まだ入れていない。2ループ用のツールは、Kira 3.1、FIRE 7.1、AMFlow 2.0、pySecDec 1.6.6、feyntrop 1.1.0、TSIL 1.46 を入れた（第~\ref{sec:twoloop} 節）。熱的な道具は、レビューに挙がったものをすべて入れ、壁の速さの WallGo も加えた（第~\ref{sec:thermal} 節）。WallGo で平衡からずれた粒子の摩擦を求めるには、WallGoMatrix と WallGoCollision が要る。メモリ 16 GB の制約から、大規模な IBP 簡約はこの Mac では難しい。退避した自作の模型を git で管理するかどうかも、まだ決めていない。

\appendix

\section{ファイルの配置}\label{app:layout}

導入後の主なファイルの配置を示す。ソースとビルドのログは \pth{hep/src} に残してある。

\begin{Verbatim}
~/hep/
  README.md                     runbook (English)
  bin/   form tform lt LoopTools fcc f++ mcc mg5_aMC wolframscript(-> Wolfram 2.app) wopen
  lib/   libooptools.a libcollier.dylib cmake/
  include/  looptools.h clooptools.h collier.mod
  mathematica/  FeynArts-3.12/  FormCalc-9.10/ (MacOSX-ARM64/: ReadForm, form 4.2.1, ...)
  MG5_aMC_v3_8_1/  (HEPTools/: ninja, oneloop, collier 1.2.9)
  opt/Wolfram142 -> /Applications/Wolfram 2.app
  src/   tarballs, build trees, *.log
~/Library/Wolfram/Applications/
  FeynCalc/  (FeynArts/ = patched 3.12, AddOns/FeynHelpers/,
              AddOns/FeynHelpers/ExternalTools/LoopTools/LoopTools -> ~/hep/bin/LoopTools)
  FeynRules/ (FeynRules.m, Core/Loop.m patched; *.orig kept)
  X/         (Package-X source, PVReduce.m)
  FeynArts.m FormCalc.m   loader stubs -> ~/hep/mathematica
  GroupMath/  Matchete/   thermal tools (2026-10-11)
~/Library/Wolfram/Paclets/Repository/
  DRalgo__DRalgo-1.5.0/  FindBounce-1.1.0/  PT2GW-1.1.0/   thermal tools (paclets)
~/miniforge3/envs/qft/      Python 3.12 for MadGraph; cosmoTransitions 2.0.7, WallGo 1.1.2 (pip), h5py
~/hep/PhaseTracer-2.2.2/  ~/hep/BSMPT-3.2.0/  ~/hep/TransitionListener-2.1.0/   thermal tools (C++, C++, Python)
~/hep/lib/libalglib.a  ~/hep/include/alglib/                 ALGLIB 4.08.0 for PhaseTracer2
~/miniforge3/envs/tl/       Python 3.12 for TransitionListener (ultranest, mpi4py, getdist, ...)
~/hep/Kira-3.1/  ~/hep/fermat-7.9b/  ~/hep/FIRE-7.1 -> src/fire-7.1/  ~/hep/AMFlow-2.0/  ~/hep/feyntrop-1.1.0/
                            two-loop tools (2026-10-12); ~/hep/bin/kira sets FERMATPATH
~/hep/lib/  libcln.a libginac.a libtsil.a libyaml-cpp.a   ~/hep/bin/  ginsh mpsolve kira
~/hep/opt/buildtools/  (meson, ninja)   ~/hep/opt/llvmwrap/  (clang wrappers)   ~/hep/opt/gmpinc/
~/Documents/Claude/QFTcalculation/
  environment/  environment.tex  files/  tests/
  models/backup_20261009/   user models (README inside)
\end{Verbatim}

\section{検算に使った値の導出}\label{app:derivations}

\paragraph{$B_0(2;1,1)$。}LoopTools の規格化では、2点関数は
\begin{equation}
\begin{split}
B_0(s;m_0,m_1)&=\frac{(2\pi\mu)^{4-D}}{i\pi^2}\int d^Dq\,\frac{1}{(q^2-m_0^2)\bigl((q+p)^2-m_1^2\bigr)}\\
&=\Delta-\int_0^1 dx\,\ln\frac{x m_1^2+(1-x)m_0^2-x(1-x)s-i\varepsilon}{\mu^2}
\end{split}
\end{equation}
である~\cite{fcalc}。ここで $s=p^2$、$\Delta=2/(4-D)-\gamma_E+\ln4\pi$ であり、LoopTools の既定値は $\Delta=0$、$\mu=1$ である。等質量 $m_0=m_1=m$ では、$a=s/m^2$ と置いて
\begin{equation}
B_0(s;m,m)=\Delta-\ln\frac{m^2}{\mu^2}-I(a),\qquad I(a)=\int_0^1dx\,\ln D(x),\quad D(x)=ax^2-ax+1
\end{equation}
となる。$0<a<4$ では $D(x)>0$ なので $i\varepsilon$ は効かない。$I(a)$ を部分積分すると、$D(1)=1$ より境界項は消え、
\begin{equation}
I(a)=\bigl[x\ln D\bigr]_0^1-\int_0^1dx\,\frac{xD'(x)}{D(x)}=-\int_0^1dx\,\frac{2ax^2-ax}{D(x)}.
\end{equation}
分子を $2ax^2-ax=2D(x)+(ax-2)$ と分け、さらに $ax-2=\tfrac12D'(x)+\tfrac{a-4}{2}$ と書くと、$\int_0^1D'/D\,dx=[\ln D]_0^1=0$ より
\begin{equation}
I(a)=-2-\frac{a-4}{2}\int_0^1\frac{dx}{D(x)}.
\end{equation}
$D(x)=a\bigl[(x-\tfrac12)^2+b^2\bigr]$、$b^2=(4-a)/(4a)>0$ と平方完成すると、
\begin{equation}
\int_0^1\frac{dx}{D(x)}=\frac{1}{ab}\Bigl[\arctan\frac{x-1/2}{b}\Bigr]_0^1=\frac{2}{ab}\arctan\frac{1}{2b}.
\end{equation}
$\tilde\beta=\sqrt{4/a-1}$ と置くと、$(4-a)/(ab)=2\tilde\beta$、$1/(2b)=1/\tilde\beta$ なので
\begin{equation}
I(a)=-2+2\tilde\beta\arctan\frac{1}{\tilde\beta},\qquad
B_0(s;m,m)=\Delta-\ln\frac{m^2}{\mu^2}+2-2\tilde\beta\arctan\frac{1}{\tilde\beta}.
\end{equation}
$s=2$、$m=\mu=1$、$\Delta=0$ では $a=2$、$\tilde\beta=1$、$\arctan1=\pi/4$ だから、$B_0(2;1,1)=2-\pi/2=0.42920367320510338\ldots$ である。

\paragraph{Package-X の $\mathrm{DiscB}$。}FeynHelpers が展開した Package-X の式では
\begin{equation}
\mathrm{DiscB}(s;m,m)=\frac{\sqrt{s(s-4m^2)}}{s}\ln\frac{2m^2-s+\sqrt{s(s-4m^2)}}{2m^2}
\end{equation}
である。$s=2$、$m=1$ では $\sqrt{s(s-4m^2)}=\sqrt{-4}=2i$、対数の引数は $(2-2+2i)/2=i$ で $\ln i=i\pi/2$ だから、$\mathrm{DiscB}=2i\cdot(i\pi/2)/2=-\pi/2$ となる。よって Package-X の有限部分 $2+\mathrm{DiscB}+\ln(\mu^2/m^2)$ は $2-\pi/2$ で、上の結果と一致する。FeynHelpers が返した値は $2-\pi/2-\gamma_E-\ln\pi=0.4292036732-0.5772156649-1.1447298858=-1.2927418775$ である。

\paragraph{$C_0(0,0,0;1,1,1)$。}外運動量がすべてゼロなら3本のプロパゲータはどれも $q^2-m^2$ で、積分は有限（$D=4$）である。Wick 回転 $q^0=iq_E^0$（$d^4q=i\,d^4q_E$、$q^2=-q_E^2$）と $\int d^4q_E=\pi^2\int_0^\infty t\,dt$（$t=q_E^2$）を使うと、
\begin{equation}
C_0=\frac{1}{i\pi^2}\int\frac{d^4q}{(q^2-m^2)^3}=-\frac{1}{\pi^2}\int\frac{d^4q_E}{(q_E^2+m^2)^3}=-\int_0^\infty\frac{t\,dt}{(t+m^2)^3}.
\end{equation}
$t/(t+m^2)^3=1/(t+m^2)^2-m^2/(t+m^2)^3$ より、最後の積分は $1/m^2-m^2/(2m^4)=1/(2m^2)$ である。よって $C_0=-1/(2m^2)$ で、$m=1$ では $-1/2$ となる。

\paragraph{$e^+e^-\to\mu^+\mu^-$ の断面積。}出発点は、質量のないレプトンに対する最低次の $\gamma$ と $Z$ 交換の微分断面積である。
\begin{equation}
\frac{d\sigma}{d\cos\theta}=\frac{\pi\alpha^2}{2s}\Bigl[(1+\cos^2\theta)A_0+2\cos\theta\,A_1\Bigr],
\end{equation}
\begin{equation}
A_0=1+2v^2\,\mathrm{Re}\,\chi+(v^2+a^2)^2|\chi|^2,\quad
A_1=2a^2\,\mathrm{Re}\,\chi+4v^2a^2|\chi|^2,\quad
\chi=\frac{1}{4s_W^2c_W^2}\,\frac{s}{s-M_Z^2+iM_Z\Gamma_Z}.
\end{equation}
ここで電子とミューオンの結合は等しく、$v=-\tfrac12+2s_W^2$、$a=-\tfrac12$ である（$T_3=-\tfrac12$、$Q=-1$）。$\chi=0$ とすると、全断面積は純 QED の $4\pi\alpha^2/(3s)$ に戻る。MadGraph の \texttt{sm} の電弱パラメータは $\alpha^{-1}(M_Z)=132.507$、$G_F=1.16639\times10^{-5}$~GeV$^{-2}$、$M_Z=91.188$~GeV、$\Gamma_Z=2.441404$~GeV である。$M_W$ と $s_W^2$ は
\begin{equation}
M_W^2=\frac{M_Z^2}{2}+\sqrt{\frac{M_Z^4}{4}-\frac{\pi\alpha M_Z^2}{\sqrt2\,G_F}},\qquad s_W^2=1-\frac{M_W^2}{M_Z^2}
\end{equation}
から、$M_W=80.4190$~GeV、$s_W^2=0.22225$ と決まる。重心系で二つのミューオンは背中合わせなので、$|\eta_\ell|<2.5$ は $|\cos\theta|<c_0=\tanh2.5=0.98661$ と同値である。$p_T>10$~GeV は $\sin\theta>0.02$、すなわち $|\cos\theta|<0.9998$ でこれより弱い。$\Delta R_{\ell\ell}>0.4$ は背中合わせなので常に満たされる。積分範囲が対称なので $\cos\theta$ に比例する項は消え、
\begin{equation}
\sigma=\frac{\pi\alpha^2}{2s}A_0\int_{-c_0}^{c_0}(1+c^2)\,dc=\frac{\pi\alpha^2}{2s}A_0\Bigl(2c_0+\frac{2c_0^3}{3}\Bigr).
\end{equation}
$\sqrt s=1$~TeV では $A_0=1.14522$ となる。$\pi\alpha^2/(2s)=8.9465\times10^{-11}$~GeV$^{-2}$、$2c_0+2c_0^3/3=2.61347$、換算係数 $0.3893794\times10^9$~pb\,GeV$^2$ を掛けて、$\sigma=0.10426$~pb を得る。この計算は \nolinkurl{tests/ee_mumu_analytic.py} で再現できる。

\section{変更・作成したファイルの全文}\label{app:files}

いずれも \nolinkurl{environment/files/} に同じものを置いてある。

\paragraph{\texttt{\textasciitilde/.zshrc} に追加したブロック。}
\VerbatimInput{files/zshrc_block.sh}

\paragraph{代替の \texttt{mcc}（\texttt{\textasciitilde/hep/bin/mcc}）。}
\VerbatimInput{files/mcc}

\paragraph{MadGraph の起動スクリプト（\texttt{\textasciitilde/hep/bin/mg5\_aMC}）。}
\VerbatimInput{files/mg5_aMC}

\paragraph{FeynRules.m への局所パッチ（\texttt{diff -u}）。}
\VerbatimInput{files/FeynRules.m.patch}

\paragraph{FeynRules の \texttt{Core/Loop.m}（NLOCT）への局所パッチ（\texttt{diff -u}）。}
\VerbatimInput{files/Loop.m.patch}

\paragraph{ノートブックを開くスクリプト（\texttt{\textasciitilde/hep/bin/wopen}）。}
\VerbatimInput{files/wopen}

\paragraph{Kira の起動スクリプト（\texttt{\textasciitilde/hep/bin/kira}）。}
\VerbatimInput{files/kira}

\paragraph{Homebrew の LLVM を使うためのスクリプト（\texttt{\textasciitilde/hep/opt/llvmwrap/}）。}\texttt{clang}、\texttt{clang++}、FIRE 用の \texttt{clang++-fire} の順に示す。
\VerbatimInput{files/llvmwrap-clang}
\VerbatimInput{files/llvmwrap-clang++}
\VerbatimInput{files/llvmwrap-clang++-fire}

\paragraph{Applications の入口ファイル。}
\VerbatimInput{files/FeynArts.m}
\VerbatimInput{files/FormCalc.m}

\paragraph{FeynHelpers 用のリンク。}
\begin{Verbatim}
FH=~/Library/Wolfram/Applications/FeynCalc/AddOns/FeynHelpers
mkdir -p $FH/ExternalTools/LoopTools
ln -sf ~/hep/bin/LoopTools $FH/ExternalTools/LoopTools/LoopTools
\end{Verbatim}

ビルドと導入に使ったコマンドの記録は \nolinkurl{files/build_commands.txt} にある。

\begin{thebibliography}{99}
\small
\bibitem{fc102} V. Shtabovenko, ``FeynCalc 10.2 and FeynHelpers 2: Multiloop calculations streamlined,'' \href{https://arxiv.org/abs/2512.19858}{arXiv:2512.19858}.
\bibitem{fc10} V. Shtabovenko, R. Mertig, F. Orellana, ``FeynCalc 10: Do multiloop integrals dream of computer codes?,'' Comput.Phys.Commun. \textbf{306} (2025) 109357, \href{https://arxiv.org/abs/2312.14089}{arXiv:2312.14089}.
\bibitem{fh1} V. Shtabovenko, ``FeynHelpers: Connecting FeynCalc to FIRE and Package-X,'' Comput.Phys.Commun. \textbf{218} (2017) 48, \href{https://arxiv.org/abs/1611.06793}{arXiv:1611.06793}.
\bibitem{fa3} T. Hahn, ``Generating Feynman diagrams and amplitudes with FeynArts 3,'' Comput.Phys.Commun. \textbf{140} (2001) 418, \href{https://arxiv.org/abs/hep-ph/0012260}{arXiv:hep-ph/0012260}.
\bibitem{fcalc} T. Hahn, M. Perez-Victoria, ``Automatized one loop calculations in four-dimensions and D-dimensions,'' Comput.Phys.Commun. \textbf{118} (1999) 153, \href{https://arxiv.org/abs/hep-ph/9807565}{arXiv:hep-ph/9807565}.
\bibitem{collier} A. Denner, S. Dittmaier, L. Hofer, ``COLLIER: a fortran-based Complex One-Loop LIbrary in Extended Regularizations,'' Comput.Phys.Commun. \textbf{212} (2017) 220, \href{https://arxiv.org/abs/1604.06792}{arXiv:1604.06792}.
\bibitem{px1} H. H. Patel, ``Package-X: A Mathematica package for the analytic calculation of one-loop integrals,'' Comput.Phys.Commun. \textbf{197} (2015) 276, \href{https://arxiv.org/abs/1503.01469}{arXiv:1503.01469}.
\bibitem{px2} H. H. Patel, ``Package-X 2.0: A Mathematica package for the analytic calculation of one-loop integrals,'' Comput.Phys.Commun. \textbf{218} (2017) 66, \href{https://arxiv.org/abs/1612.00009}{arXiv:1612.00009}.
\bibitem{form5} J. Davies \textit{et al.}, ``FORM Version 5.0,'' \href{https://arxiv.org/abs/2601.19982}{arXiv:2601.19982}.
\bibitem{fr2} A. Alloul \textit{et al.}, ``FeynRules 2.0 -- A complete toolbox for tree-level phenomenology,'' Comput.Phys.Commun. \textbf{185} (2014) 2250, \href{https://arxiv.org/abs/1310.1921}{arXiv:1310.1921}.
\bibitem{nloct} C. Degrande, ``Automatic evaluation of UV and R2 terms for beyond the Standard Model Lagrangians: a proof-of-principle,'' Comput.Phys.Commun. \textbf{197} (2015) 239, \href{https://arxiv.org/abs/1406.3030}{arXiv:1406.3030}.
\bibitem{mg5} J. Alwall \textit{et al.}, ``The automated computation of tree-level and next-to-leading order differential cross sections, and their matching to parton shower simulations,'' JHEP \textbf{07} (2014) 079, \href{https://arxiv.org/abs/1405.0301}{arXiv:1405.0301}.
\bibitem{ninja} T. Peraro, ``Ninja: Automated Integrand Reduction via Laurent Expansion for One-Loop Amplitudes,'' Comput.Phys.Commun. \textbf{185} (2014) 2771, \href{https://arxiv.org/abs/1403.1229}{arXiv:1403.1229}.
\bibitem{ninjamg} V. Hirschi, T. Peraro, ``Tensor integrand reduction via Laurent expansion,'' JHEP \textbf{06} (2016) 060, \href{https://arxiv.org/abs/1604.01363}{arXiv:1604.01363}.
\bibitem{oneloop} A. van Hameren, ``OneLOop: For the evaluation of one-loop scalar functions,'' Comput.Phys.Commun. \textbf{182} (2011) 2427, \href{https://arxiv.org/abs/1007.4716}{arXiv:1007.4716}.
\bibitem{allanach} B. Allanach, ``LHC di-lepton searches for $Z'$ bosons which explain measurements of $b\to s\ell^+\ell^-$ transitions,'' Phys.Lett.B \textbf{858} (2024) 139055, \href{https://arxiv.org/abs/2404.14748}{arXiv:2404.14748}（付属 UFO は FeynRules 2.3.49 と Mathematica 14.0.0 で生成）.
\bibitem{url:form} FORM releases. \url{https://github.com/form-dev/form/releases}
\bibitem{url:feynarts} FeynArts, FormCalc, LoopTools. \url{https://feynarts.de/}
\bibitem{url:collier} COLLIER. \url{https://collier.hepforge.org/}
\bibitem{url:px} Package-X（McMule による再公開）. \url{https://mule-tools.gitlab.io/package-x/}
\bibitem{url:fr} FeynRules. \url{https://feynrules.irmp.ucl.ac.be/}
\bibitem{url:mg5} MadGraph5\_aMC@NLO. \url{https://launchpad.net/mg5amcnlo}
\bibitem{url:feynmaster} D. Fontes, J. C. Rom\~ao, FeynMaster Manual v2（Mathematica 14 での FeynRules の \texttt{MatrixSymbol} の扱い）. \url{https://porthos.tecnico.ulisboa.pt/FeynMaster/Manual/Manual-v2.pdf}
\bibitem{phi34thermal} $\phi^3+\phi^4$ 理論の有限温度の一次相転移と重力波（本リポジトリ \nolinkurl{models/phi34_thermal/ResearchNote/note.pdf}）。熱的な道具との照合は第7節と付録~H。
\bibitem{cosmotransitions} C. L. Wainwright, ``CosmoTransitions: computing cosmological phase transition temperatures and bubble profiles with multiple fields,'' Comput.Phys.Commun. \textbf{183} (2012) 2006, \href{https://arxiv.org/abs/1109.4189}{arXiv:1109.4189}.
\bibitem{dralgo} A. Ekstedt, P. Schicho, T. V. I. Tenkanen, ``DRalgo: a package for effective field theory approach for thermal phase transitions,'' Comput.Phys.Commun. \textbf{288} (2023) 108725, \href{https://arxiv.org/abs/2205.08815}{arXiv:2205.08815}.
\bibitem{groupmath} R. M. Fonseca, ``GroupMath: a Mathematica package for group theory calculations,'' Comput.Phys.Commun. \textbf{267} (2021) 108085, \href{https://arxiv.org/abs/2011.01764}{arXiv:2011.01764}.
\bibitem{matchete} J. Fuentes-Mart\'{\i}n, M. K\"onig, J. Pag\`es, A. E. Thomsen, F. Wilsch, ``A proof of concept for matchete: an automated tool for matching effective theories,'' Eur.Phys.J.C \textbf{83} (2023) 662, \href{https://arxiv.org/abs/2212.04510}{arXiv:2212.04510}.
\bibitem{matchotter} J. Fuentes-Mart\'{\i}n, J. L\'opez Miras, A. Moreno-S\'anchez, ``Matchotter: an automated tool for dimensional reduction at finite temperature,'' \href{https://arxiv.org/abs/2604.21972}{arXiv:2604.21972}.
\bibitem{findbounce} V. Guada, M. Nemev\v{s}ek, M. Pintar, ``FindBounce: package for multi-field bounce actions,'' Comput.Phys.Commun. \textbf{256} (2020) 107480, \href{https://arxiv.org/abs/2002.00881}{arXiv:2002.00881}.
\bibitem{pt2gw} V. Brdar, M. Finetti, M. Matteini, A. P. Morais, M. Nemev\v{s}ek, ``PT2GWFinder: a package for cosmological first-order phase transitions and gravitational waves,'' Comput.Phys.Commun. \textbf{323} (2026) 110119, \href{https://arxiv.org/abs/2505.04744}{arXiv:2505.04744}.
\bibitem{caprini24} C. Caprini \textit{et al.}, ``Gravitational waves from first-order phase transitions in LISA: reconstruction pipeline and physics interpretation,'' JCAP \textbf{10} (2024) 020, \href{https://arxiv.org/abs/2403.03723}{arXiv:2403.03723}.
\bibitem{url:dralgo} DRalgo. \url{https://github.com/DR-algo/DRalgo}
\bibitem{url:groupmath} GroupMath. \url{https://renatofonseca.net/groupmath}
\bibitem{url:matchete} Matchete. \url{https://gitlab.com/matchete/matchete}
\bibitem{url:findbounce} FindBounce. \url{https://github.com/vguada/FindBounce}
\bibitem{url:pt2gw} PT2GWFinder. \url{https://github.com/finshky/PT2GW}
\bibitem{url:ct} CosmoTransitions（PyPI）. \url{https://pypi.org/project/cosmoTransitions/}
\bibitem{phasetracer2} P. Athron \textit{et al.}, ``PhaseTracer2: from the effective potential to gravitational waves,'' Eur.Phys.J.C \textbf{85} (2025) 559, \href{https://arxiv.org/abs/2412.04881}{arXiv:2412.04881}.
\bibitem{bsmpt3} P. Basler \textit{et al.}, ``BSMPT v3: a tool for phase transitions and primordial gravitational waves in extended Higgs sectors,'' \href{https://arxiv.org/abs/2404.19037}{arXiv:2404.19037}.
\bibitem{tl2} J. Matuszak, C. Tasillo, ``TransitionListener v2.0 -- robust gravitational wave predictions for cosmological phase transitions,'' \href{https://arxiv.org/abs/2605.15259}{arXiv:2605.15259}.
\bibitem{url:phasetracer} PhaseTracer. \url{https://github.com/PhaseTracer/PhaseTracer}
\bibitem{url:alglib} ALGLIB. \url{https://www.alglib.net/}
\bibitem{url:bsmpt} BSMPT. \url{https://github.com/BSMPT/BSMPT}
\bibitem{url:tl} TransitionListener. \url{https://github.com/tasicarl/TransitionListener}
\bibitem{wallgo} A. Ekstedt, O. Gould, J. Hirvonen, B. Laurent, L. Niemi, P. Schicho, J. van de Vis, ``How fast does the WallGo? A package for computing wall velocities in first-order phase transitions,'' JHEP \textbf{04} (2025) 101, \href{https://arxiv.org/abs/2411.04970}{arXiv:2411.04970}.
\bibitem{url:wallgo} WallGo（PyPI）. \url{https://pypi.org/project/WallGo/}
\bibitem{phi3note} $\phi^3$ 理論：模型の定義と1ループの構造（本リポジトリ \nolinkurl{models/phi3/ResearchNote/note.pdf}）。2ループの道具との照合は第8節と付録~T、U。
\bibitem{kira3} F. Lange, J. Usovitsch, Z. Wu, ``Kira 3: integral reduction with efficient seeding and optimized equation selection,'' Comput.Phys.Commun. \textbf{322} (2026) 109999, \href{https://arxiv.org/abs/2505.20197}{arXiv:2505.20197}.
\bibitem{fire7} A. V. Smirnov, M. Zeng, ``FIRE 7: Automatic Reduction with Modular Approach,'' Comput.Phys.Commun. \textbf{325} (2026) 110200, \href{https://arxiv.org/abs/2510.07150}{arXiv:2510.07150}.
\bibitem{amflow} X. Liu, Y.-Q. Ma, ``AMFlow: a Mathematica package for Feynman integrals computation via auxiliary mass flow,'' Comput.Phys.Commun. \textbf{283} (2023) 108565, \href{https://arxiv.org/abs/2201.11669}{arXiv:2201.11669}.
\bibitem{amflow2} R.-J. Huang, X. Liu, Y.-Q. Ma, ``AMFlow 2.0: significant algorithmic and software improvements for Feynman integral evaluation,'' \href{https://arxiv.org/abs/2607.08477}{arXiv:2607.08477}.
\bibitem{pysecdec} S. Borowka \textit{et al.}, ``pySecDec: a toolbox for the numerical evaluation of multi-scale integrals,'' Comput.Phys.Commun. \textbf{222} (2018) 313, \href{https://arxiv.org/abs/1703.09692}{arXiv:1703.09692}.
\bibitem{feyntrop} M. Borinsky, H. J. Munch, F. Tellander, ``Tropical Feynman integration in the Minkowski regime,'' Comput.Phys.Commun. \textbf{292} (2023) 108874, \href{https://arxiv.org/abs/2302.08955}{arXiv:2302.08955}.
\bibitem{tsil} S. P. Martin, D. G. Robertson, ``TSIL: a program for the calculation of two-loop self-energy integrals,'' Comput.Phys.Commun. \textbf{174} (2006) 133, \href{https://arxiv.org/abs/hep-ph/0501132}{hep-ph/0501132}.
\bibitem{url:kira} Kira. \url{https://gitlab.com/kira-pyred/kira}
\bibitem{url:fermat} Fermat（R. H. Lewis）. \url{http://home.bway.net/lewis/}
\bibitem{url:ginac} GiNaC、CLN. \url{https://www.ginac.de/}
\bibitem{url:fire} FIRE. \url{https://gitlab.srcc.msu.ru/feynmanintegrals/fire}
\bibitem{url:amflow} AMFlow. \url{https://gitlab.com/multiloop-pku/amflow}
\bibitem{url:pysecdec} pySecDec（PyPI）. \url{https://pypi.org/project/pySecDec/}
\bibitem{url:feyntrop} feyntrop. \url{https://github.com/michibo/feyntrop}
\bibitem{url:tsil} TSIL. \url{https://www.niu.edu/spmartin/TSIL/}
\end{thebibliography}

\end{document}
